% Generated by scripts/build_interaction_cases.py; do not edit payloads.
\section{Full Interactions for the Discussion Cases}
\label{app:discussion-interactions}
This section reproduces the complete exported shared conversations behind
Figure~\ref{fig:discussion-cases}, including every recorded message and
handoff field. C2 documents workflow revision; C1 documents clarification
with persona 155 and its matched no-persona interaction. All are original
trial-1 endpoints using GPT-6 Astra with low reasoning, and each received
a normalized score of 100/100. Persona 155's full configured profile
appears in Appendix~\ref{app:persona-profile}.

The original Chinese and English wording is retained without abridgment.
Payload formatting, mathematical source notation and structured fields are
printed literally; line wrapping is adjusted for the page. The translated
main-text excerpts use editorial highlights and omission marks. Source
event numbers are preserved, including gaps corresponding to activity
outside the shared conversation. Private tool logs, system prompts and
evaluator material are outside this export. Source paths, event-log hashes
and byte-identical conversation snapshots are retained in the paper's
\texttt{review/discussion\_candidates} archive.

The workflow transcript refers to limma~\cite{ritchie2015limma},
Benjamini--Hochberg (BH) correction~\cite{benjamini1995fdr}, and
Holm correction~\cite{holm1979sequential}.

% The bundled OFL font is a renamed, character-subsetted Noto Sans CJK SC.
% XeLaTeX/Tectonic is required for original-language conversation text.
\font\interactionfont="[fonts/PSBTraceCJK.otf]" at 8.5pt
\newenvironment{interactionevent}[1]{%
  \begin{tcolorbox}[breakable,colback=gray!2,colframe=gray!45,
    boxrule=0.4pt,arc=1pt,left=7pt,right=7pt,top=4pt,bottom=4pt,
    before skip=5pt,after skip=5pt,lines before break=4,
    title={#1},title after break={#1 (continued)},
    colbacktitle=gray!12,coltitle=black,fonttitle=\small\bfseries]
  \interactionfont\baselineskip=11.5pt
  \raggedright
  \parindent=0pt\parskip=1.5pt\emergencystretch=1em
}{\end{tcolorbox}}
\newcommand{\interactionfield}[1]{\par{\normalfont\small\bfseries #1}\par\nobreak}

\Needspace{10\baselineskip}
\subsection{Workflow Revision: BixBench with Persona 155}
\label{app:case-workflow}
Task: \texttt{bix-\allowbreak{}3-\allowbreak{}q5}. Trial 1; original endpoint; persona 155.
% Snapshot SHA256: 22f9264e725a7ba2efe70457a145df2ebb20a8fafeba737de4faf5f83cb48f97
\begin{interactionevent}{Event 1: Researcher (simulator)}
\interactionfield{Message}
请\allowbreak{}先\allowbreak{}只\allowbreak{}检\allowbreak{}查\allowbreak{}输\allowbreak{}入\allowbreak{}文\allowbreak{}件\allowbreak{}与\allowbreak{}现\allowbreak{}有\allowbreak{}分\allowbreak{}析\allowbreak{}资\allowbreak{}料，\allowbreak{}不\allowbreak{}要\allowbreak{}开\allowbreak{}始\allowbreak{}完\allowbreak{}整\allowbreak{}分\allowbreak{}析。\allowbreak{}请\allowbreak{}核\allowbreak{}查\allowbreak{} data/\allowbreak{}Data\_\allowbreak{}deposition\_\allowbreak{}RNAseq\_\allowbreak{}Paroxetine\_\allowbreak{}2017.xlsx 两\allowbreak{}个\allowbreak{}工\allowbreak{}作\allowbreak{}表\allowbreak{}的\allowbreak{}实\allowbreak{}际\allowbreak{}结\allowbreak{}构、\allowbreak{}基\allowbreak{}因\allowbreak{}标\allowbreak{}识、\allowbreak{}表\allowbreak{}达\allowbreak{}值\allowbreak{}性\allowbreak{}质、\allowbreak{}各\allowbreak{}组\allowbreak{}织\allowbreak{}样\allowbreak{}本\allowbreak{}数，\allowbreak{}以\allowbreak{}及\allowbreak{}样\allowbreak{}本\allowbreak{}编\allowbreak{}号\allowbreak{}与\allowbreak{}元\allowbreak{}数\allowbreak{}据\allowbreak{}能\allowbreak{}否\allowbreak{}对\allowbreak{}应；\allowbreak{}特\allowbreak{}别\allowbreak{}说\allowbreak{}明\allowbreak{}基\allowbreak{}线\allowbreak{}血\allowbreak{}液\allowbreak{}与\allowbreak{}末\allowbreak{}次\allowbreak{}血\allowbreak{}液\allowbreak{}是\allowbreak{}否\allowbreak{}存\allowbreak{}在\allowbreak{}可\allowbreak{}确\allowbreak{}认\allowbreak{}的\allowbreak{}配\allowbreak{}对\allowbreak{}关\allowbreak{}系。\allowbreak{}随\allowbreak{}后\allowbreak{}提\allowbreak{}出\allowbreak{}三\allowbreak{}项\allowbreak{}比\allowbreak{}较\allowbreak{}（Final vs Baseline Blood、\allowbreak{}Dentate gyrus vs Baseline blood、\allowbreak{}Dentate gyrus vs Final blood）的\allowbreak{}差\allowbreak{}异\allowbreak{}表\allowbreak{}达\allowbreak{}分\allowbreak{}析\allowbreak{}方\allowbreak{}案，\allowbreak{}并\allowbreak{}说\allowbreak{}明\allowbreak{}精\allowbreak{}确\allowbreak{}二\allowbreak{}项\allowbreak{}检\allowbreak{}验\allowbreak{}的\allowbreak{}试\allowbreak{}验\allowbreak{}单\allowbreak{}位、\allowbreak{}成\allowbreak{}功\allowbreak{}事\allowbreak{}件、\allowbreak{}零\allowbreak{}假\allowbreak{}设\allowbreak{}概\allowbreak{}率\allowbreak{}及\allowbreak{}5\%阈\allowbreak{}值\allowbreak{}的\allowbreak{}用\allowbreak{}途。\allowbreak{}请\allowbreak{}区\allowbreak{}分\allowbreak{}文\allowbreak{}件\allowbreak{}明\allowbreak{}确\allowbreak{}规\allowbreak{}定\allowbreak{}的\allowbreak{}方\allowbreak{}法\allowbreak{}与\allowbreak{}您\allowbreak{}的\allowbreak{}建\allowbreak{}议，\allowbreak{}列\allowbreak{}出\allowbreak{}会\allowbreak{}影\allowbreak{}响\allowbreak{}结\allowbreak{}论\allowbreak{}的\allowbreak{}缺\allowbreak{}失\allowbreak{}信\allowbreak{}息、\allowbreak{}调\allowbreak{}整\allowbreak{}因\allowbreak{}素\allowbreak{}和\allowbreak{}假\allowbreak{}设。\allowbreak{}此\allowbreak{}次\allowbreak{}仅\allowbreak{}提\allowbreak{}交\allowbreak{}检\allowbreak{}查\allowbreak{}结\allowbreak{}果\allowbreak{}与\allowbreak{}拟\allowbreak{}议\allowbreak{}方\allowbreak{}案。\allowbreak{}\par
\interactionfield{Interaction}
initial\_\allowbreak{}request\par
\interactionfield{Phase}
initial\_\allowbreak{}inspection\par
\interactionfield{Action}
continue\par
\interactionfield{Approve Approach}
false\par
\interactionfield{Evidence}
已\allowbreak{}提\allowbreak{}供\allowbreak{}的\allowbreak{}结\allowbreak{}构\allowbreak{}信\allowbreak{}息\allowbreak{}显\allowbreak{}示\allowbreak{}工\allowbreak{}作\allowbreak{}簿\allowbreak{}包\allowbreak{}含\allowbreak{} Mapping\_\allowbreak{}statistics 和\allowbreak{} NormCount，\allowbreak{}后\allowbreak{}者\allowbreak{}的\allowbreak{}首\allowbreak{}行\allowbreak{}标\allowbreak{}有\allowbreak{}三\allowbreak{}类\allowbreak{}组\allowbreak{}织；\allowbreak{}尚\allowbreak{}未\allowbreak{}核\allowbreak{}查\allowbreak{}具\allowbreak{}体\allowbreak{}样\allowbreak{}本、\allowbreak{}表\allowbreak{}达\allowbreak{}值\allowbreak{}及\allowbreak{}检\allowbreak{}验\allowbreak{}定\allowbreak{}义，\allowbreak{}不\allowbreak{}能\allowbreak{}据\allowbreak{}此\allowbreak{}判\allowbreak{}断\allowbreak{}哪\allowbreak{}项\allowbreak{}比\allowbreak{}较\allowbreak{}效\allowbreak{}应\allowbreak{}可\allowbreak{}忽\allowbreak{}略。\allowbreak{}\par
\interactionfield{Objectives Open}
[\par
  "核\allowbreak{}实\allowbreak{}输\allowbreak{}入\allowbreak{}结\allowbreak{}构、\allowbreak{}样\allowbreak{}本\allowbreak{}分\allowbreak{}组\allowbreak{}及\allowbreak{}配\allowbreak{}对\allowbreak{}关\allowbreak{}系\allowbreak{}",\par
  "明\allowbreak{}确\allowbreak{}差\allowbreak{}异\allowbreak{}表\allowbreak{}达\allowbreak{}方\allowbreak{}法\allowbreak{}与\allowbreak{}精\allowbreak{}确\allowbreak{}二\allowbreak{}项\allowbreak{}检\allowbreak{}验\allowbreak{}的\allowbreak{}定\allowbreak{}义\allowbreak{}及\allowbreak{}5\%阈\allowbreak{}值\allowbreak{}",\par
  "审\allowbreak{}议\allowbreak{}方\allowbreak{}案\allowbreak{}后\allowbreak{}执\allowbreak{}行\allowbreak{}三\allowbreak{}项\allowbreak{}比\allowbreak{}较\allowbreak{}并\allowbreak{}确\allowbreak{}定\allowbreak{}效\allowbreak{}应\allowbreak{}可\allowbreak{}忽\allowbreak{}略\allowbreak{}的\allowbreak{}比\allowbreak{}较\allowbreak{}",\par
  "保\allowbreak{}留\allowbreak{}已\allowbreak{}执\allowbreak{}行\allowbreak{}分\allowbreak{}析\allowbreak{}脚\allowbreak{}本\allowbreak{}或\allowbreak{}笔\allowbreak{}记\allowbreak{}本\allowbreak{}路\allowbreak{}径\allowbreak{}并\allowbreak{}提\allowbreak{}供\allowbreak{}简\allowbreak{}洁\allowbreak{}结\allowbreak{}论\allowbreak{}"\par
]\par
\end{interactionevent}
\begin{interactionevent}{Event 12: Actor}
\interactionfield{Answer}
\#\# 检\allowbreak{}查\allowbreak{}结\allowbreak{}果\allowbreak{}（尚\allowbreak{}未\allowbreak{}执\allowbreak{}行\allowbreak{}差\allowbreak{}异\allowbreak{}表\allowbreak{}达\allowbreak{}或\allowbreak{}二\allowbreak{}项\allowbreak{}检\allowbreak{}验）\par
工\allowbreak{}作\allowbreak{}目\allowbreak{}录\allowbreak{}仅\allowbreak{}有\allowbreak{} task.json 与\allowbreak{}该\allowbreak{} Excel，\allowbreak{}无\allowbreak{}现\allowbreak{}成\allowbreak{}脚\allowbreak{}本、\allowbreak{}论\allowbreak{}文\allowbreak{}方\allowbreak{}法\allowbreak{}或\allowbreak{}比\allowbreak{}较\allowbreak{}编\allowbreak{}号\allowbreak{}说\allowbreak{}明。\allowbreak{}\par
\smallskip
**Mapping\_\allowbreak{}statistics：\allowbreak{}91行\allowbreak{}×22列。\allowbreak{}** 第\allowbreak{}1行\allowbreak{}为\allowbreak{}表\allowbreak{}头，\allowbreak{}随\allowbreak{}后\allowbreak{}90个\allowbreak{}样\allowbreak{}本。\allowbreak{}有\allowbreak{}效\allowbreak{}核\allowbreak{}心\allowbreak{}字\allowbreak{}段\allowbreak{}为\allowbreak{} Index、\allowbreak{}Sample ID、\allowbreak{}total read、\allowbreak{}NH:i:1、\allowbreak{}HI:I:0、\allowbreak{}mapping rate、\allowbreak{}Response、\allowbreak{}Tissue、\allowbreak{}Code、\allowbreak{}DESeq（htseq\_\allowbreak{}count 文\allowbreak{}件\allowbreak{}名）、\allowbreak{}Index\_\allowbreak{}DESeq。\allowbreak{}第\allowbreak{}12–16列\allowbreak{}主\allowbreak{}要\allowbreak{}是\allowbreak{}用\allowbreak{}于\allowbreak{}拼\allowbreak{}接\allowbreak{}代\allowbreak{}码\allowbreak{}的\allowbreak{}引\allowbreak{}号\allowbreak{}及\allowbreak{}组\allowbreak{}别\allowbreak{}字\allowbreak{}符\allowbreak{}串，\allowbreak{}并\allowbreak{}非\allowbreak{}表\allowbreak{}头\allowbreak{}所\allowbreak{}称\allowbreak{}的\allowbreak{}额\allowbreak{}外\allowbreak{}比\allowbreak{}对\allowbreak{}统\allowbreak{}计；\allowbreak{}第\allowbreak{}17–22列\allowbreak{}数\allowbreak{}据\allowbreak{}为\allowbreak{}空。\allowbreak{}不\allowbreak{}能\allowbreak{}把\allowbreak{}这\allowbreak{}些\allowbreak{}重\allowbreak{}复\allowbreak{}表\allowbreak{}头\allowbreak{}当\allowbreak{}成\allowbreak{}独\allowbreak{}立\allowbreak{}测\allowbreak{}序\allowbreak{}批\allowbreak{}次。\allowbreak{}\par
\smallskip
**NormCount：\allowbreak{}25,406行\allowbreak{}×188列。\allowbreak{}** 第\allowbreak{}1–3行\allowbreak{}为\allowbreak{}分\allowbreak{}组\allowbreak{}/\allowbreak{}说\allowbreak{}明，\allowbreak{}第\allowbreak{}4行\allowbreak{}为\allowbreak{}字\allowbreak{}段\allowbreak{}名，\allowbreak{}第\allowbreak{}5行\allowbreak{}起\allowbreak{}25,402个\allowbreak{}基\allowbreak{}因。\allowbreak{}第\allowbreak{}1–2列\allowbreak{}为\allowbreak{} Index/\allowbreak{}GeneID；\allowbreak{}第\allowbreak{}3–92列\allowbreak{}为\allowbreak{}90个\allowbreak{}样\allowbreak{}本；\allowbreak{}第\allowbreak{}93–98列\allowbreak{}为\allowbreak{}三\allowbreak{}个\allowbreak{}组\allowbreak{}织\allowbreak{}的\allowbreak{} Max/\allowbreak{}Median；\allowbreak{}第\allowbreak{}99–179列\allowbreak{}为\allowbreak{}9组\allowbreak{}既\allowbreak{}有\allowbreak{}“DESeq, Comparison N, exclude outliers”结\allowbreak{}果，\allowbreak{}每\allowbreak{}组\allowbreak{}含\allowbreak{} baseMean、\allowbreak{}baseMeanA/\allowbreak{}B、\allowbreak{}foldChange、\allowbreak{}log2FoldChange、\allowbreak{}pval、\allowbreak{}padj、\allowbreak{}两\allowbreak{}列\allowbreak{}DEG；\allowbreak{}第\allowbreak{}180–188列\allowbreak{}为\allowbreak{}另\allowbreak{}9列\allowbreak{}DEG摘\allowbreak{}要。\allowbreak{}**这\allowbreak{}9个\allowbreak{}编\allowbreak{}号\allowbreak{}没\allowbreak{}有\allowbreak{}可\allowbreak{}核\allowbreak{}实\allowbreak{}的\allowbreak{}对\allowbreak{}比\allowbreak{}定\allowbreak{}义，\allowbreak{}不\allowbreak{}能\allowbreak{}直\allowbreak{}接\allowbreak{}当\allowbreak{}成\allowbreak{}本\allowbreak{}题\allowbreak{}三\allowbreak{}项\allowbreak{}组\allowbreak{}织\allowbreak{}对\allowbreak{}比。\allowbreak{}** 文\allowbreak{}件\allowbreak{}未\allowbreak{}给\allowbreak{}出\allowbreak{}明\allowbreak{}确\allowbreak{}异\allowbreak{}常\allowbreak{}样\allowbreak{}本\allowbreak{}清\allowbreak{}单、\allowbreak{}DEG阈\allowbreak{}值\allowbreak{}或\allowbreak{}模\allowbreak{}型\allowbreak{}设\allowbreak{}计。\allowbreak{}公\allowbreak{}式\allowbreak{}检\allowbreak{}查\allowbreak{}见\allowbreak{}到\allowbreak{}的\allowbreak{}是\allowbreak{}“-\allowbreak{}Inf”记\allowbreak{}录，\allowbreak{}不\allowbreak{}是\allowbreak{}可\allowbreak{}复\allowbreak{}现\allowbreak{}的\allowbreak{}DEG判\allowbreak{}定\allowbreak{}规\allowbreak{}则。\allowbreak{}\par
\smallskip
GeneID为\allowbreak{}25,402个\allowbreak{}非\allowbreak{}空、\allowbreak{}唯\allowbreak{}一\allowbreak{}的\allowbreak{}符\allowbreak{}号\allowbreak{}式\allowbreak{}标\allowbreak{}识，\allowbreak{}例\allowbreak{}如\allowbreak{}0610007C21Rik、\allowbreak{}Zyx、\allowbreak{}Zzz3；\allowbreak{}格\allowbreak{}式\allowbreak{}符\allowbreak{}合\allowbreak{}小\allowbreak{}鼠\allowbreak{}基\allowbreak{}因\allowbreak{}符\allowbreak{}号，\allowbreak{}但\allowbreak{}文\allowbreak{}件\allowbreak{}没\allowbreak{}有\allowbreak{}提\allowbreak{}供\allowbreak{}物\allowbreak{}种\allowbreak{}/\allowbreak{}注\allowbreak{}释\allowbreak{}版\allowbreak{}本\allowbreak{}的\allowbreak{}正\allowbreak{}式\allowbreak{}说\allowbreak{}明，\allowbreak{}也\allowbreak{}不\allowbreak{}是\allowbreak{}Ensembl稳\allowbreak{}定\allowbreak{}ID。\allowbreak{}表\allowbreak{}达\allowbreak{}矩\allowbreak{}阵\allowbreak{}无\allowbreak{}缺\allowbreak{}失、\allowbreak{}非\allowbreak{}负，\allowbreak{}范\allowbreak{}围\allowbreak{}0–57,771,177.4；\allowbreak{}2,286,180个\allowbreak{}值\allowbreak{}中\allowbreak{}1,571,838个\allowbreak{}非\allowbreak{}整\allowbreak{}数，\allowbreak{}2,210个\allowbreak{}基\allowbreak{}因\allowbreak{}全\allowbreak{}零。\allowbreak{}表\allowbreak{}名\allowbreak{}及\allowbreak{}数\allowbreak{}值\allowbreak{}支\allowbreak{}持\allowbreak{}“归\allowbreak{}一\allowbreak{}化\allowbreak{}计\allowbreak{}数\allowbreak{}”，\allowbreak{}**不\allowbreak{}是\allowbreak{}可\allowbreak{}直\allowbreak{}接\allowbreak{}交\allowbreak{}给\allowbreak{}DESeq2的\allowbreak{}原\allowbreak{}始\allowbreak{}整\allowbreak{}数\allowbreak{}计\allowbreak{}数\allowbreak{}**；\allowbreak{}具\allowbreak{}体\allowbreak{}归\allowbreak{}一\allowbreak{}化\allowbreak{}方\allowbreak{}法\allowbreak{}及\allowbreak{}是\allowbreak{}否\allowbreak{}跨\allowbreak{}组\allowbreak{}织\allowbreak{}统\allowbreak{}一\allowbreak{}标\allowbreak{}度\allowbreak{}未\allowbreak{}说\allowbreak{}明，\allowbreak{}不\allowbreak{}能\allowbreak{}擅\allowbreak{}称\allowbreak{}TPM/\allowbreak{}FPKM或\allowbreak{}取\allowbreak{}整\allowbreak{}恢\allowbreak{}复\allowbreak{}原\allowbreak{}始\allowbreak{}计\allowbreak{}数。\allowbreak{}\par
\smallskip
|组\allowbreak{}织\allowbreak{}|Control|good\_\allowbreak{}responder|bad\_\allowbreak{}responder|合\allowbreak{}计\allowbreak{}|\par
|-\allowbreak{}-\allowbreak{}-\allowbreak{}|-\allowbreak{}-\allowbreak{}-:|-\allowbreak{}-\allowbreak{}-:|-\allowbreak{}-\allowbreak{}-:|-\allowbreak{}-\allowbreak{}-:|\par
|baseline\_\allowbreak{}blood|10|10|10|30|\par
|dentate\_\allowbreak{}gyrus|10|10|10|30|\par
|final\_\allowbreak{}blood|10|10|10|30|\par
\smallskip
90个\allowbreak{}表\allowbreak{}达\allowbreak{}列\allowbreak{}名\allowbreak{}与\allowbreak{}元\allowbreak{}数\allowbreak{}据\allowbreak{}Code完\allowbreak{}整\allowbreak{}一\allowbreak{}对\allowbreak{}一\allowbreak{}匹\allowbreak{}配，\allowbreak{}Index\_\allowbreak{}DESeq也\allowbreak{}与\allowbreak{}表\allowbreak{}达\allowbreak{}列\allowbreak{}顺\allowbreak{}序\allowbreak{}逐\allowbreak{}一\allowbreak{}一\allowbreak{}致。\allowbreak{}\par
\smallskip
\#\#\# 配\allowbreak{}对\allowbreak{}证\allowbreak{}据\allowbreak{}\par
**基\allowbreak{}线\allowbreak{}与\allowbreak{}末\allowbreak{}次\allowbreak{}血\allowbreak{}液\allowbreak{}具\allowbreak{}有\allowbreak{}完\allowbreak{}整、\allowbreak{}可\allowbreak{}由\allowbreak{}文\allowbreak{}件\allowbreak{}编\allowbreak{}号\allowbreak{}核\allowbreak{}查\allowbreak{}的\allowbreak{}一\allowbreak{}对\allowbreak{}一\allowbreak{}对\allowbreak{}应，\allowbreak{}不\allowbreak{}应\allowbreak{}仅\allowbreak{}因\allowbreak{}没\allowbreak{}有\allowbreak{}Subject列\allowbreak{}就\allowbreak{}按\allowbreak{}独\allowbreak{}立\allowbreak{}样\allowbreak{}本\allowbreak{}处\allowbreak{}理。\allowbreak{}** Sample ID末\allowbreak{}尾\allowbreak{}组\allowbreak{}织\allowbreak{}后\allowbreak{}缀\allowbreak{}之\allowbreak{}前\allowbreak{}的\allowbreak{}数\allowbreak{}字，\allowbreak{}在\allowbreak{}三\allowbreak{}个\allowbreak{}组\allowbreak{}织\allowbreak{}中\allowbreak{}重\allowbreak{}复\allowbreak{}且\allowbreak{}组\allowbreak{}别\allowbreak{}一\allowbreak{}致：\allowbreak{}例\allowbreak{}如\allowbreak{}Control的\allowbreak{}98对\allowbreak{}应\allowbreak{}baseline `...\_\allowbreak{}08\_\allowbreak{}98\_\allowbreak{}bl`、\allowbreak{}final `...\_\allowbreak{}09\_\allowbreak{}98\_\allowbreak{}fi`、\allowbreak{}齿\allowbreak{}状\allowbreak{}回\allowbreak{} `...\_\allowbreak{}07\_\allowbreak{}98\_\allowbreak{}c`。\allowbreak{}全\allowbreak{}部\allowbreak{}30个\allowbreak{}编\allowbreak{}号\allowbreak{}在\allowbreak{}各\allowbreak{}组\allowbreak{}织\allowbreak{}各\allowbreak{}出\allowbreak{}现\allowbreak{}一\allowbreak{}次，\allowbreak{}且\allowbreak{}与\allowbreak{}各\allowbreak{}Response组\allowbreak{}的\allowbreak{}Code尾\allowbreak{}号\allowbreak{}对\allowbreak{}应\allowbreak{}一\allowbreak{}致。\allowbreak{}因\allowbreak{}此\allowbreak{}可\allowbreak{}建\allowbreak{}立\allowbreak{}30对\allowbreak{}血\allowbreak{}样，\allowbreak{}也\allowbreak{}可\allowbreak{}建\allowbreak{}立\allowbreak{}30个\allowbreak{}三\allowbreak{}组\allowbreak{}织\allowbreak{}样\allowbreak{}本\allowbreak{}组。\allowbreak{}严\allowbreak{}谨\allowbreak{}界\allowbreak{}限\allowbreak{}是：\allowbreak{}文\allowbreak{}件\allowbreak{}未\allowbreak{}明\allowbreak{}确\allowbreak{}声\allowbreak{}明\allowbreak{}该\allowbreak{}数\allowbreak{}字\allowbreak{}为\allowbreak{}动\allowbreak{}物\allowbreak{}ID；\allowbreak{}“同\allowbreak{}一\allowbreak{}动\allowbreak{}物\allowbreak{}重\allowbreak{}复\allowbreak{}测\allowbreak{}量\allowbreak{}”仍\allowbreak{}是\allowbreak{}有\allowbreak{}强\allowbreak{}编\allowbreak{}号\allowbreak{}证\allowbreak{}据\allowbreak{}支\allowbreak{}持\allowbreak{}的\allowbreak{}假\allowbreak{}设，\allowbreak{}而\allowbreak{}非\allowbreak{}独\allowbreak{}立\allowbreak{}实\allowbreak{}验\allowbreak{}记\allowbreak{}录\allowbreak{}的\allowbreak{}确\allowbreak{}认。\allowbreak{}\par
另\allowbreak{}有\allowbreak{}一\allowbreak{}项\allowbreak{}命\allowbreak{}名\allowbreak{}异\allowbreak{}常：\allowbreak{}`bad\_\allowbreak{}responder\_\allowbreak{}dentate\_\allowbreak{}gyrus2` 的\allowbreak{}Sample ID为\allowbreak{} `...\_\allowbreak{}23\_\allowbreak{}102\_\allowbreak{}bl`，\allowbreak{}后\allowbreak{}缀\allowbreak{}像\allowbreak{}基\allowbreak{}线\allowbreak{}血；\allowbreak{}Tissue与\allowbreak{}Code均\allowbreak{}标\allowbreak{}为\allowbreak{}齿\allowbreak{}状\allowbreak{}回。\allowbreak{}拟\allowbreak{}以\allowbreak{}显\allowbreak{}式\allowbreak{}Tissue/\allowbreak{}Code为\allowbreak{}主\allowbreak{}并\allowbreak{}标\allowbreak{}记\allowbreak{}异\allowbreak{}常，\allowbreak{}做\allowbreak{}排\allowbreak{}除\allowbreak{}此\allowbreak{}编\allowbreak{}号\allowbreak{}的\allowbreak{}敏\allowbreak{}感\allowbreak{}性\allowbreak{}检\allowbreak{}查，\allowbreak{}不\allowbreak{}静\allowbreak{}默\allowbreak{}改\allowbreak{}组\allowbreak{}织。\allowbreak{}\par
\smallskip
\#\# 拟\allowbreak{}议\allowbreak{}分\allowbreak{}析\allowbreak{}方\allowbreak{}案\allowbreak{}\par
**任\allowbreak{}务\allowbreak{}明\allowbreak{}确\allowbreak{}要\allowbreak{}求：\allowbreak{}** 三\allowbreak{}项\allowbreak{}指\allowbreak{}定\allowbreak{}比\allowbreak{}较，\allowbreak{}以\allowbreak{}及\allowbreak{}精\allowbreak{}确\allowbreak{}二\allowbreak{}项\allowbreak{}检\allowbreak{}验\allowbreak{}和\allowbreak{}5\%阈\allowbreak{}值。\allowbreak{}**文\allowbreak{}件\allowbreak{}明\allowbreak{}确\allowbreak{}记\allowbreak{}录：\allowbreak{}** 有\allowbreak{}归\allowbreak{}一\allowbreak{}化\allowbreak{}计\allowbreak{}数\allowbreak{}及\allowbreak{}既\allowbreak{}有\allowbreak{}DESeq、\allowbreak{}exclude outliers结\allowbreak{}果。\allowbreak{}文\allowbreak{}件\allowbreak{}没\allowbreak{}有\allowbreak{}明\allowbreak{}确\allowbreak{}要\allowbreak{}求\allowbreak{}新\allowbreak{}分\allowbreak{}析\allowbreak{}使\allowbreak{}用\allowbreak{}哪\allowbreak{}一\allowbreak{}版\allowbreak{}DESeq、\allowbreak{}哪\allowbreak{}些\allowbreak{}协\allowbreak{}变\allowbreak{}量、\allowbreak{}何\allowbreak{}种\allowbreak{}过\allowbreak{}滤\allowbreak{}或\allowbreak{}二\allowbreak{}项\allowbreak{}检\allowbreak{}验\allowbreak{}定\allowbreak{}义。\allowbreak{}以\allowbreak{}下\allowbreak{}均\allowbreak{}为\allowbreak{}我\allowbreak{}的\allowbreak{}预\allowbreak{}先\allowbreak{}建\allowbreak{}议，\allowbreak{}而\allowbreak{}非\allowbreak{}文\allowbreak{}件\allowbreak{}规\allowbreak{}定。\allowbreak{}\par
\smallskip
1. 依\allowbreak{}Code连\allowbreak{}接\allowbreak{}元\allowbreak{}数\allowbreak{}据，\allowbreak{}按\allowbreak{}上\allowbreak{}述\allowbreak{}数\allowbreak{}字\allowbreak{}建\allowbreak{}立\allowbreak{}候\allowbreak{}选\allowbreak{}Subject；\allowbreak{}保\allowbreak{}留\allowbreak{}三\allowbreak{}个\allowbreak{}Response层\allowbreak{}各\allowbreak{}10个\allowbreak{}对\allowbreak{}象。\allowbreak{}预\allowbreak{}先\allowbreak{}过\allowbreak{}滤\allowbreak{}低\allowbreak{}表\allowbreak{}达\allowbreak{}基\allowbreak{}因，\allowbreak{}例\allowbreak{}如\allowbreak{}至\allowbreak{}少\allowbreak{}10个\allowbreak{}样\allowbreak{}本\allowbreak{}归\allowbreak{}一\allowbreak{}化\allowbreak{}计\allowbreak{}数\allowbreak{}≥10，\allowbreak{}三\allowbreak{}项\allowbreak{}比\allowbreak{}较\allowbreak{}使\allowbreak{}用\allowbreak{}同\allowbreak{}一\allowbreak{}个\allowbreak{}过\allowbreak{}滤\allowbreak{}基\allowbreak{}因\allowbreak{}集\allowbreak{}合，\allowbreak{}并\allowbreak{}报\allowbreak{}告\allowbreak{}过\allowbreak{}滤\allowbreak{}及\allowbreak{}可\allowbreak{}检\allowbreak{}验\allowbreak{}数\allowbreak{}量。\allowbreak{}该\allowbreak{}阈\allowbreak{}值\allowbreak{}为\allowbreak{}建\allowbreak{}议，\allowbreak{}需\allowbreak{}敏\allowbreak{}感\allowbreak{}性\allowbreak{}检\allowbreak{}查。\allowbreak{}\par
2. 在\allowbreak{}只\allowbreak{}有\allowbreak{}归\allowbreak{}一\allowbreak{}化\allowbreak{}计\allowbreak{}数\allowbreak{}的\allowbreak{}情\allowbreak{}况\allowbreak{}下，\allowbreak{}建\allowbreak{}议\allowbreak{}以\allowbreak{}log2(count+1)进\allowbreak{}行\allowbreak{}limma经\allowbreak{}验\allowbreak{}贝\allowbreak{}叶\allowbreak{}斯\allowbreak{}线\allowbreak{}性\allowbreak{}模\allowbreak{}型，\allowbreak{}以\allowbreak{}Subject固\allowbreak{}定\allowbreak{}效\allowbreak{}应\allowbreak{}控\allowbreak{}制\allowbreak{}重\allowbreak{}复\allowbreak{}测\allowbreak{}量，\allowbreak{}组\allowbreak{}织\allowbreak{}为\allowbreak{}目\allowbreak{}标\allowbreak{}因\allowbreak{}素；\allowbreak{}不\allowbreak{}要\allowbreak{}把\allowbreak{}这\allowbreak{}些\allowbreak{}小\allowbreak{}数\allowbreak{}当\allowbreak{}原\allowbreak{}始\allowbreak{}计\allowbreak{}数\allowbreak{}运\allowbreak{}行\allowbreak{}DESeq2，\allowbreak{}也\allowbreak{}不\allowbreak{}把\allowbreak{}它\allowbreak{}们\allowbreak{}作\allowbreak{}为\allowbreak{}原\allowbreak{}始\allowbreak{}计\allowbreak{}数\allowbreak{}做\allowbreak{}voom。\allowbreak{}三\allowbreak{}个\allowbreak{}方\allowbreak{}向\allowbreak{}固\allowbreak{}定\allowbreak{}为\allowbreak{}Final−Baseline、\allowbreak{}Dentate−Baseline、\allowbreak{}Dentate−Final，\allowbreak{}报\allowbreak{}告\allowbreak{}log2尺\allowbreak{}度\allowbreak{}差\allowbreak{}异、\allowbreak{}原\allowbreak{}始\allowbreak{}p及\allowbreak{}每\allowbreak{}项\allowbreak{}比\allowbreak{}较\allowbreak{}内\allowbreak{}BH校\allowbreak{}正\allowbreak{}p。\allowbreak{}此\allowbreak{}方\allowbreak{}案\allowbreak{}是\allowbreak{}对\allowbreak{}现\allowbreak{}有\allowbreak{}数\allowbreak{}据\allowbreak{}的\allowbreak{}近\allowbreak{}似\allowbreak{}分\allowbreak{}析，\allowbreak{}不\allowbreak{}是\allowbreak{}重\allowbreak{}现\allowbreak{}原\allowbreak{}始\allowbreak{}DESeq流\allowbreak{}程。\allowbreak{}若\allowbreak{}取\allowbreak{}得\allowbreak{}原\allowbreak{}始\allowbreak{}计\allowbreak{}数\allowbreak{}及\allowbreak{}设\allowbreak{}计\allowbreak{}说\allowbreak{}明，\allowbreak{}应\allowbreak{}优\allowbreak{}先\allowbreak{}按\allowbreak{}说\allowbreak{}明\allowbreak{}用\allowbreak{}计\allowbreak{}数\allowbreak{}模\allowbreak{}型\allowbreak{}重\allowbreak{}分\allowbreak{}析。\allowbreak{}\par
3. Subject效\allowbreak{}应\allowbreak{}已\allowbreak{}吸\allowbreak{}收\allowbreak{}不\allowbreak{}随\allowbreak{}组\allowbreak{}织\allowbreak{}改\allowbreak{}变\allowbreak{}的\allowbreak{}Response主\allowbreak{}效\allowbreak{}应，\allowbreak{}不\allowbreak{}能\allowbreak{}再\allowbreak{}重\allowbreak{}复\allowbreak{}加\allowbreak{}入\allowbreak{}造\allowbreak{}成\allowbreak{}秩\allowbreak{}亏。\allowbreak{}主\allowbreak{}结\allowbreak{}果\allowbreak{}为\allowbreak{}三\allowbreak{}个\allowbreak{}Response层\allowbreak{}等\allowbreak{}权\allowbreak{}平\allowbreak{}均\allowbreak{}的\allowbreak{}组\allowbreak{}织\allowbreak{}/\allowbreak{}时\allowbreak{}间\allowbreak{}差\allowbreak{}异；\allowbreak{}加\allowbreak{}入\allowbreak{}Response×组\allowbreak{}织\allowbreak{}交\allowbreak{}互\allowbreak{}或\allowbreak{}分\allowbreak{}层\allowbreak{}结\allowbreak{}果\allowbreak{}作\allowbreak{}敏\allowbreak{}感\allowbreak{}性\allowbreak{}分\allowbreak{}析，\allowbreak{}避\allowbreak{}免\allowbreak{}把\allowbreak{}处\allowbreak{}理\allowbreak{}反\allowbreak{}应\allowbreak{}异\allowbreak{}质\allowbreak{}性\allowbreak{}混\allowbreak{}为\allowbreak{}纯\allowbreak{}组\allowbreak{}织\allowbreak{}效\allowbreak{}应。\allowbreak{}未\allowbreak{}经\allowbreak{}确\allowbreak{}证\allowbreak{}不\allowbreak{}按\allowbreak{}表\allowbreak{}达\allowbreak{}差\allowbreak{}异\allowbreak{}挑\allowbreak{}选\allowbreak{}异\allowbreak{}常\allowbreak{}样\allowbreak{}本；\allowbreak{}先\allowbreak{}做\allowbreak{}分\allowbreak{}布\allowbreak{}/\allowbreak{}PCA等\allowbreak{}QC并\allowbreak{}报\allowbreak{}告。\allowbreak{}候\allowbreak{}选\allowbreak{}配\allowbreak{}对\allowbreak{}与\allowbreak{}不\allowbreak{}配\allowbreak{}对\allowbreak{}模\allowbreak{}型\allowbreak{}也\allowbreak{}可\allowbreak{}对\allowbreak{}照，\allowbreak{}但\allowbreak{}后\allowbreak{}者\allowbreak{}不\allowbreak{}是\allowbreak{}主\allowbreak{}要\allowbreak{}设\allowbreak{}计。\allowbreak{}\par
\smallskip
\#\#\# 精\allowbreak{}确\allowbreak{}二\allowbreak{}项\allowbreak{}检\allowbreak{}验\allowbreak{}的\allowbreak{}明\allowbreak{}确\allowbreak{}操\allowbreak{}作\allowbreak{}定\allowbreak{}义\allowbreak{}（建\allowbreak{}议）\par
-\allowbreak{} **试\allowbreak{}验\allowbreak{}单\allowbreak{}位：\allowbreak{}** 每\allowbreak{}项\allowbreak{}比\allowbreak{}较\allowbreak{}中\allowbreak{}一\allowbreak{}个\allowbreak{}通\allowbreak{}过\allowbreak{}预\allowbreak{}过\allowbreak{}滤\allowbreak{}且\allowbreak{}得\allowbreak{}到\allowbreak{}有\allowbreak{}效\allowbreak{}检\allowbreak{}验\allowbreak{}p值\allowbreak{}的\allowbreak{}基\allowbreak{}因，\allowbreak{}而\allowbreak{}非\allowbreak{}动\allowbreak{}物、\allowbreak{}样\allowbreak{}本、\allowbreak{}reads或\allowbreak{}三\allowbreak{}个\allowbreak{}比\allowbreak{}较\allowbreak{}本\allowbreak{}身。\allowbreak{}\par
-\allowbreak{} **成\allowbreak{}功\allowbreak{}事\allowbreak{}件：\allowbreak{}** 该\allowbreak{}基\allowbreak{}因\allowbreak{}双\allowbreak{}侧\allowbreak{}差\allowbreak{}异\allowbreak{}表\allowbreak{}达\allowbreak{}检\allowbreak{}验\allowbreak{}的\allowbreak{}**未\allowbreak{}校\allowbreak{}正\allowbreak{}p<0.05**。\allowbreak{}令\allowbreak{}成\allowbreak{}功\allowbreak{}数\allowbreak{}k、\allowbreak{}有\allowbreak{}效\allowbreak{}基\allowbreak{}因\allowbreak{}数\allowbreak{}n。\allowbreak{}\par
-\allowbreak{} **零\allowbreak{}假\allowbreak{}设\allowbreak{}概\allowbreak{}率：\allowbreak{}** p\ensuremath{{}_0}=\allowbreak{}0.05，\allowbreak{}即\allowbreak{}在\allowbreak{}全\allowbreak{}局\allowbreak{}无\allowbreak{}差\allowbreak{}异\allowbreak{}且\allowbreak{}基\allowbreak{}因\allowbreak{}检\allowbreak{}验\allowbreak{}校\allowbreak{}准\allowbreak{}成\allowbreak{}立\allowbreak{}时，\allowbreak{}单\allowbreak{}基\allowbreak{}因\allowbreak{}名\allowbreak{}义\allowbreak{}显\allowbreak{}著\allowbreak{}的\allowbreak{}概\allowbreak{}率\allowbreak{}为\allowbreak{}5\%；\allowbreak{}拟\allowbreak{}检\allowbreak{}验\allowbreak{}Binomial(n,0.05)的\allowbreak{}上\allowbreak{}尾，\allowbreak{}判\allowbreak{}断\allowbreak{}显\allowbreak{}著\allowbreak{}基\allowbreak{}因\allowbreak{}是\allowbreak{}否\allowbreak{}超\allowbreak{}出\allowbreak{}偶\allowbreak{}然\allowbreak{}预\allowbreak{}期。\allowbreak{}\par
-\allowbreak{} **5\%的\allowbreak{}用\allowbreak{}途\allowbreak{}分\allowbreak{}开：\allowbreak{}** ①单\allowbreak{}基\allowbreak{}因\allowbreak{}名\allowbreak{}义\allowbreak{}筛\allowbreak{}选\allowbreak{}阈\allowbreak{}值\allowbreak{}0.05，\allowbreak{}因\allowbreak{}而\allowbreak{}定\allowbreak{}义\allowbreak{}p\ensuremath{{}_0}；\allowbreak{}②二\allowbreak{}项\allowbreak{}检\allowbreak{}验\allowbreak{}显\allowbreak{}著\allowbreak{}性\allowbreak{}水\allowbreak{}平\allowbreak{}0.05。\allowbreak{}另\allowbreak{}报\allowbreak{}告\allowbreak{}BH-\allowbreak{}FDR<0.05的\allowbreak{}DEG作\allowbreak{}为\allowbreak{}多\allowbreak{}重\allowbreak{}校\allowbreak{}正\allowbreak{}结\allowbreak{}果，\allowbreak{}但\allowbreak{}**不\allowbreak{}能\allowbreak{}将\allowbreak{}BH显\allowbreak{}著\allowbreak{}数\allowbreak{}代\allowbreak{}入\allowbreak{}后\allowbreak{}仍\allowbreak{}声\allowbreak{}称\allowbreak{}其\allowbreak{}零\allowbreak{}假\allowbreak{}设\allowbreak{}成\allowbreak{}功\allowbreak{}率\allowbreak{}精\allowbreak{}确\allowbreak{}等\allowbreak{}于\allowbreak{}0.05**。\allowbreak{}三\allowbreak{}项\allowbreak{}二\allowbreak{}项\allowbreak{}检\allowbreak{}验\allowbreak{}报\allowbreak{}告\allowbreak{}原\allowbreak{}始\allowbreak{}p及\allowbreak{}Holm校\allowbreak{}正\allowbreak{}p。\allowbreak{}\par
-\allowbreak{} 文\allowbreak{}件\allowbreak{}并\allowbreak{}未\allowbreak{}规\allowbreak{}定\allowbreak{}上\allowbreak{}述\allowbreak{}二\allowbreak{}项\allowbreak{}检\allowbreak{}验\allowbreak{}细\allowbreak{}节，\allowbreak{}特\allowbreak{}别\allowbreak{}是\allowbreak{}没\allowbreak{}有\allowbreak{}规\allowbreak{}定\allowbreak{}“5\%”为\allowbreak{}可\allowbreak{}忽\allowbreak{}略\allowbreak{}效\allowbreak{}应\allowbreak{}比\allowbreak{}例\allowbreak{}上\allowbreak{}限。\allowbreak{}上\allowbreak{}尾\allowbreak{}不\allowbreak{}显\allowbreak{}著\allowbreak{}只\allowbreak{}表\allowbreak{}示\allowbreak{}未\allowbreak{}检\allowbreak{}出\allowbreak{}超\allowbreak{}过\allowbreak{}偶\allowbreak{}然\allowbreak{}水\allowbreak{}平\allowbreak{}的\allowbreak{}显\allowbreak{}著\allowbreak{}基\allowbreak{}因，\allowbreak{}**不\allowbreak{}证\allowbreak{}明\allowbreak{}生\allowbreak{}物\allowbreak{}学\allowbreak{}等\allowbreak{}效\allowbreak{}或\allowbreak{}无\allowbreak{}效\allowbreak{}应\allowbreak{}**。\allowbreak{}若\allowbreak{}要\allowbreak{}严\allowbreak{}格\allowbreak{}证\allowbreak{}明\allowbreak{}效\allowbreak{}应\allowbreak{}可\allowbreak{}忽\allowbreak{}略，\allowbreak{}需\allowbreak{}要\allowbreak{}预\allowbreak{}先\allowbreak{}指\allowbreak{}定\allowbreak{}效\allowbreak{}应\allowbreak{}大\allowbreak{}小\allowbreak{}/\allowbreak{}比\allowbreak{}例\allowbreak{}界\allowbreak{}限\allowbreak{}并\allowbreak{}做\allowbreak{}等\allowbreak{}效\allowbreak{}性\allowbreak{}分\allowbreak{}析，\allowbreak{}不\allowbreak{}能\allowbreak{}事\allowbreak{}后\allowbreak{}把\allowbreak{}不\allowbreak{}显\allowbreak{}著\allowbreak{}改\allowbreak{}称\allowbreak{}等\allowbreak{}效。\allowbreak{}\par
-\allowbreak{} 基\allowbreak{}因\allowbreak{}间\allowbreak{}相\allowbreak{}关\allowbreak{}违\allowbreak{}反\allowbreak{}二\allowbreak{}项\allowbreak{}独\allowbreak{}立\allowbreak{}性，\allowbreak{}故\allowbreak{}按\allowbreak{}要\allowbreak{}求\allowbreak{}计\allowbreak{}算\allowbreak{}的\allowbreak{}“精\allowbreak{}确\allowbreak{}”p仅\allowbreak{}在\allowbreak{}二\allowbreak{}项\allowbreak{}模\allowbreak{}型\allowbreak{}假\allowbreak{}设\allowbreak{}下\allowbreak{}精\allowbreak{}确；\allowbreak{}结\allowbreak{}合\allowbreak{}k/\allowbreak{}n、\allowbreak{}效\allowbreak{}应\allowbreak{}分\allowbreak{}布\allowbreak{}和\allowbreak{}敏\allowbreak{}感\allowbreak{}性\allowbreak{}检\allowbreak{}查\allowbreak{}解\allowbreak{}释，\allowbreak{}不\allowbreak{}单\allowbreak{}凭\allowbreak{}极\allowbreak{}小\allowbreak{}p排\allowbreak{}序。\allowbreak{}\par
\smallskip
\#\# 会\allowbreak{}影\allowbreak{}响\allowbreak{}结\allowbreak{}论\allowbreak{}的\allowbreak{}缺\allowbreak{}失\allowbreak{}信\allowbreak{}息\allowbreak{}与\allowbreak{}假\allowbreak{}设\allowbreak{}\par
关\allowbreak{}键\allowbreak{}缺\allowbreak{}失\allowbreak{}为\allowbreak{}原\allowbreak{}始\allowbreak{}计\allowbreak{}数、\allowbreak{}归\allowbreak{}一\allowbreak{}化\allowbreak{}方\allowbreak{}法\allowbreak{}及\allowbreak{}跨\allowbreak{}组\allowbreak{}织\allowbreak{}可\allowbreak{}比\allowbreak{}性、\allowbreak{}动\allowbreak{}物\allowbreak{}ID编\allowbreak{}码\allowbreak{}说\allowbreak{}明、\allowbreak{}9个\allowbreak{}既\allowbreak{}有\allowbreak{}比\allowbreak{}较\allowbreak{}定\allowbreak{}义、\allowbreak{}异\allowbreak{}常\allowbreak{}样\allowbreak{}本\allowbreak{}清\allowbreak{}单、\allowbreak{}处\allowbreak{}理\allowbreak{}剂\allowbreak{}量\allowbreak{}/\allowbreak{}时\allowbreak{}间、\allowbreak{}年\allowbreak{}龄\allowbreak{}性\allowbreak{}别、\allowbreak{}实\allowbreak{}验\allowbreak{}/\allowbreak{}建\allowbreak{}库\allowbreak{}批\allowbreak{}次、\allowbreak{}RNA质\allowbreak{}量\allowbreak{}及\allowbreak{}细\allowbreak{}胞\allowbreak{}组\allowbreak{}成。\allowbreak{}血\allowbreak{}液\allowbreak{}时\allowbreak{}间\allowbreak{}效\allowbreak{}应\allowbreak{}可\allowbreak{}能\allowbreak{}混\allowbreak{}有\allowbreak{}治\allowbreak{}疗、\allowbreak{}时\allowbreak{}间\allowbreak{}和\allowbreak{}细\allowbreak{}胞\allowbreak{}比\allowbreak{}例\allowbreak{}变\allowbreak{}化；\allowbreak{}脑\allowbreak{}血\allowbreak{}差\allowbreak{}异\allowbreak{}也\allowbreak{}可\allowbreak{}能\allowbreak{}混\allowbreak{}有\allowbreak{}批\allowbreak{}次\allowbreak{}和\allowbreak{}细\allowbreak{}胞\allowbreak{}组\allowbreak{}成。\allowbreak{}现\allowbreak{}有\allowbreak{}Response可\allowbreak{}以\allowbreak{}分\allowbreak{}层\allowbreak{}/\allowbreak{}交\allowbreak{}互\allowbreak{}处\allowbreak{}理，\allowbreak{}mapping rate及\allowbreak{}测\allowbreak{}序\allowbreak{}深\allowbreak{}度\allowbreak{}可\allowbreak{}作\allowbreak{}QC，\allowbreak{}但\allowbreak{}不\allowbreak{}默\allowbreak{}认\allowbreak{}把\allowbreak{}归\allowbreak{}一\allowbreak{}化\allowbreak{}后\allowbreak{}的\allowbreak{}总\allowbreak{}reads再\allowbreak{}当\allowbreak{}offset，\allowbreak{}亦\allowbreak{}不\allowbreak{}能\allowbreak{}凭\allowbreak{}空\allowbreak{}调\allowbreak{}整\allowbreak{}缺\allowbreak{}失\allowbreak{}协\allowbreak{}变\allowbreak{}量。\allowbreak{}拟\allowbreak{}采\allowbreak{}用\allowbreak{}有\allowbreak{}证\allowbreak{}据\allowbreak{}支\allowbreak{}持\allowbreak{}的\allowbreak{}配\allowbreak{}对、\allowbreak{}保\allowbreak{}持\allowbreak{}组\allowbreak{}织\allowbreak{}显\allowbreak{}式\allowbreak{}标\allowbreak{}签、\allowbreak{}把\allowbreak{}归\allowbreak{}一\allowbreak{}化\allowbreak{}值\allowbreak{}视\allowbreak{}为\allowbreak{}可\allowbreak{}比\allowbreak{}较\allowbreak{}表\allowbreak{}达\allowbreak{}尺\allowbreak{}度，\allowbreak{}并\allowbreak{}报\allowbreak{}告\allowbreak{}这\allowbreak{}些\allowbreak{}假\allowbreak{}设\allowbreak{}的\allowbreak{}限\allowbreak{}制。\allowbreak{}\par
\smallskip
请\allowbreak{}批\allowbreak{}准\allowbreak{}上\allowbreak{}述\allowbreak{}主\allowbreak{}设\allowbreak{}计\allowbreak{}及\allowbreak{}二\allowbreak{}项\allowbreak{}检\allowbreak{}验\allowbreak{}定\allowbreak{}义\allowbreak{}后\allowbreak{}再\allowbreak{}开\allowbreak{}展\allowbreak{}完\allowbreak{}整\allowbreak{}分\allowbreak{}析。\allowbreak{}本\allowbreak{}次\allowbreak{}不\allowbreak{}判\allowbreak{}断\allowbreak{}哪\allowbreak{}项\allowbreak{}比\allowbreak{}较\allowbreak{}效\allowbreak{}应\allowbreak{}可\allowbreak{}忽\allowbreak{}略，\allowbreak{}也\allowbreak{}未\allowbreak{}执\allowbreak{}行\allowbreak{}任\allowbreak{}何\allowbreak{}差\allowbreak{}异\allowbreak{}表\allowbreak{}达\allowbreak{}或\allowbreak{}二\allowbreak{}项\allowbreak{}检\allowbreak{}验。\allowbreak{}\par
\interactionfield{Stage}
plan\par
\interactionfield{Evidence}
已\allowbreak{}实\allowbreak{}际\allowbreak{}读\allowbreak{}取\allowbreak{}两\allowbreak{}个\allowbreak{}工\allowbreak{}作\allowbreak{}表\allowbreak{}全\allowbreak{}部\allowbreak{}数\allowbreak{}据\allowbreak{}用\allowbreak{}于\allowbreak{}结\allowbreak{}构\allowbreak{}检\allowbreak{}查：\allowbreak{}Mapping\_\allowbreak{}statistics 91×22，\allowbreak{}NormCount 25406×188；\allowbreak{}25402个\allowbreak{}唯\allowbreak{}一\allowbreak{}非\allowbreak{}空\allowbreak{}GeneID；\allowbreak{}90个\allowbreak{}表\allowbreak{}达\allowbreak{}列\allowbreak{}与\allowbreak{}Code及\allowbreak{}Index\_\allowbreak{}DESeq完\allowbreak{}整\allowbreak{}匹\allowbreak{}配；\allowbreak{}九\allowbreak{}个\allowbreak{}组\allowbreak{}织\allowbreak{}×Response单\allowbreak{}元\allowbreak{}各\allowbreak{}10样\allowbreak{}本；\allowbreak{}表\allowbreak{}达\allowbreak{}矩\allowbreak{}阵\allowbreak{}无\allowbreak{}缺\allowbreak{}失、\allowbreak{}2210行\allowbreak{}全\allowbreak{}零、\allowbreak{}1571838个\allowbreak{}非\allowbreak{}整\allowbreak{}数\allowbreak{}值。\allowbreak{}Sample ID包\allowbreak{}含\allowbreak{}30个\allowbreak{}跨\allowbreak{}三\allowbreak{}组\allowbreak{}织\allowbreak{}重\allowbreak{}复\allowbreak{}的\allowbreak{}候\allowbreak{}选\allowbreak{}对\allowbreak{}象\allowbreak{}数\allowbreak{}字；\allowbreak{}例\allowbreak{}98对\allowbreak{}应\allowbreak{}\_\allowbreak{}08\_\allowbreak{}98\_\allowbreak{}bl、\allowbreak{}\_\allowbreak{}09\_\allowbreak{}98\_\allowbreak{}fi、\allowbreak{}\_\allowbreak{}07\_\allowbreak{}98\_\allowbreak{}c。\allowbreak{}仅\allowbreak{}进\allowbreak{}行\allowbreak{}了\allowbreak{}输\allowbreak{}入\allowbreak{}检\allowbreak{}查，\allowbreak{}没\allowbreak{}有\allowbreak{}拟\allowbreak{}合\allowbreak{}DE模\allowbreak{}型\allowbreak{}或\allowbreak{}计\allowbreak{}算\allowbreak{}二\allowbreak{}项\allowbreak{}检\allowbreak{}验。\allowbreak{}\par
\interactionfield{Assumptions}
[\par
  "Sample ID组\allowbreak{}织\allowbreak{}后\allowbreak{}缀\allowbreak{}前\allowbreak{}的\allowbreak{}数\allowbreak{}字\allowbreak{}代\allowbreak{}表\allowbreak{}同\allowbreak{}一\allowbreak{}实\allowbreak{}验\allowbreak{}对\allowbreak{}象，\allowbreak{}拟\allowbreak{}用\allowbreak{}于\allowbreak{}配\allowbreak{}对；\allowbreak{}缺\allowbreak{}乏\allowbreak{}正\allowbreak{}式\allowbreak{}ID字\allowbreak{}典。\allowbreak{}",\par
  "NormCount为\allowbreak{}可\allowbreak{}跨\allowbreak{}样\allowbreak{}本\allowbreak{}比\allowbreak{}较\allowbreak{}的\allowbreak{}归\allowbreak{}一\allowbreak{}化\allowbreak{}计\allowbreak{}数，\allowbreak{}其\allowbreak{}确\allowbreak{}切\allowbreak{}算\allowbreak{}法\allowbreak{}仍\allowbreak{}未\allowbreak{}知。\allowbreak{}"\par
]\par
\interactionfield{Uncertainties}
[\par
  "9个\allowbreak{}既\allowbreak{}有\allowbreak{}DESeq比\allowbreak{}较\allowbreak{}的\allowbreak{}组\allowbreak{}别\allowbreak{}和\allowbreak{}异\allowbreak{}常\allowbreak{}值\allowbreak{}排\allowbreak{}除\allowbreak{}规\allowbreak{}则\allowbreak{}未\allowbreak{}说\allowbreak{}明。\allowbreak{}",\par
  "题\allowbreak{}目\allowbreak{}仅\allowbreak{}说\allowbreak{}明\allowbreak{}精\allowbreak{}确\allowbreak{}二\allowbreak{}项\allowbreak{}检\allowbreak{}验\allowbreak{}与\allowbreak{}5\%阈\allowbreak{}值，\allowbreak{}未\allowbreak{}明\allowbreak{}示\allowbreak{}成\allowbreak{}功\allowbreak{}事\allowbreak{}件、\allowbreak{}零\allowbreak{}假\allowbreak{}设\allowbreak{}概\allowbreak{}率\allowbreak{}及\allowbreak{}检\allowbreak{}验\allowbreak{}方\allowbreak{}向；\allowbreak{}本\allowbreak{}方\allowbreak{}案\allowbreak{}已\allowbreak{}将\allowbreak{}建\allowbreak{}议\allowbreak{}预\allowbreak{}先\allowbreak{}明\allowbreak{}确。\allowbreak{}",\par
  "生\allowbreak{}物\allowbreak{}学\allowbreak{}可\allowbreak{}忽\allowbreak{}略\allowbreak{}效\allowbreak{}应\allowbreak{}界\allowbreak{}限\allowbreak{}未\allowbreak{}提\allowbreak{}供，\allowbreak{}非\allowbreak{}显\allowbreak{}著\allowbreak{}不\allowbreak{}等\allowbreak{}同\allowbreak{}于\allowbreak{}等\allowbreak{}效。\allowbreak{}"\par
]\par
\end{interactionevent}
\begin{interactionevent}{Event 13: Researcher (simulator)}
\interactionfield{Message}
批\allowbreak{}准\allowbreak{}按\allowbreak{}所\allowbreak{}提\allowbreak{}方\allowbreak{}案\allowbreak{}开\allowbreak{}展\allowbreak{}实\allowbreak{}际\allowbreak{}分\allowbreak{}析。\allowbreak{}归\allowbreak{}一\allowbreak{}化\allowbreak{}小\allowbreak{}数\allowbreak{}计\allowbreak{}数\allowbreak{}不\allowbreak{}应\allowbreak{}取\allowbreak{}整\allowbreak{}后\allowbreak{}用\allowbreak{}于\allowbreak{}DESeq2；\allowbreak{}采\allowbreak{}用\allowbreak{}log2(count+1)的\allowbreak{}limma模\allowbreak{}型，\allowbreak{}建\allowbreak{}议\allowbreak{}使\allowbreak{}用\allowbreak{}trend=\allowbreak{}TRUE处\allowbreak{}理\allowbreak{}均\allowbreak{}值\allowbreak{}—方\allowbreak{}差\allowbreak{}关\allowbreak{}系，\allowbreak{}以\allowbreak{}候\allowbreak{}选\allowbreak{}Subject固\allowbreak{}定\allowbreak{}效\allowbreak{}应\allowbreak{}控\allowbreak{}制\allowbreak{}配\allowbreak{}对，\allowbreak{}完\allowbreak{}成\allowbreak{}三\allowbreak{}项\allowbreak{}指\allowbreak{}定\allowbreak{}对\allowbreak{}比。\allowbreak{}配\allowbreak{}对\allowbreak{}与\allowbreak{}跨\allowbreak{}组\allowbreak{}织\allowbreak{}归\allowbreak{}一\allowbreak{}化\allowbreak{}可\allowbreak{}比\allowbreak{}性\allowbreak{}仍\allowbreak{}须\allowbreak{}明\allowbreak{}确\allowbreak{}标\allowbreak{}为\allowbreak{}假\allowbreak{}设，\allowbreak{}不\allowbreak{}能\allowbreak{}写\allowbreak{}成\allowbreak{}已\allowbreak{}确\allowbreak{}认\allowbreak{}事\allowbreak{}实。\allowbreak{}采\allowbreak{}用\allowbreak{}预\allowbreak{}先\allowbreak{}提\allowbreak{}出\allowbreak{}的\allowbreak{}共\allowbreak{}同\allowbreak{}过\allowbreak{}滤\allowbreak{}规\allowbreak{}则\allowbreak{}（至\allowbreak{}少\allowbreak{}10个\allowbreak{}样\allowbreak{}本\allowbreak{}计\allowbreak{}数\allowbreak{}≥10），\allowbreak{}报\allowbreak{}告\allowbreak{}有\allowbreak{}效\allowbreak{}基\allowbreak{}因\allowbreak{}数、\allowbreak{}效\allowbreak{}应\allowbreak{}分\allowbreak{}布、\allowbreak{}原\allowbreak{}始\allowbreak{}p及\allowbreak{}比\allowbreak{}较\allowbreak{}内\allowbreak{}BH结\allowbreak{}果。\allowbreak{}按\allowbreak{}所\allowbreak{}提\allowbreak{}定\allowbreak{}义\allowbreak{}执\allowbreak{}行\allowbreak{}精\allowbreak{}确\allowbreak{}二\allowbreak{}项\allowbreak{}上\allowbreak{}尾\allowbreak{}检\allowbreak{}验：\allowbreak{}k为\allowbreak{}原\allowbreak{}始\allowbreak{}p<0.05的\allowbreak{}基\allowbreak{}因\allowbreak{}数，\allowbreak{}n为\allowbreak{}有\allowbreak{}效\allowbreak{}检\allowbreak{}验\allowbreak{}数，\allowbreak{}p\ensuremath{{}_0}=\allowbreak{}0.05；\allowbreak{}报\allowbreak{}告\allowbreak{}k/\allowbreak{}n、\allowbreak{}原\allowbreak{}始\allowbreak{}及\allowbreak{}Holm校\allowbreak{}正\allowbreak{}p，\allowbreak{}并\allowbreak{}说\allowbreak{}明\allowbreak{}基\allowbreak{}因\allowbreak{}相\allowbreak{}关\allowbreak{}性\allowbreak{}限\allowbreak{}制。\allowbreak{}该\allowbreak{}定\allowbreak{}义\allowbreak{}是\allowbreak{}本\allowbreak{}次\allowbreak{}操\allowbreak{}作\allowbreak{}性\allowbreak{}假\allowbreak{}设，\allowbreak{}不\allowbreak{}代\allowbreak{}表\allowbreak{}已\allowbreak{}恢\allowbreak{}复\allowbreak{}原\allowbreak{}作\allowbreak{}者\allowbreak{}方\allowbreak{}法。\allowbreak{}完\allowbreak{}成\allowbreak{}基\allowbreak{}本\allowbreak{}QC及\allowbreak{}已\allowbreak{}提\allowbreak{}出\allowbreak{}的\allowbreak{}过\allowbreak{}滤、\allowbreak{}配\allowbreak{}对、\allowbreak{}Response异\allowbreak{}质\allowbreak{}性\allowbreak{}和\allowbreak{}异\allowbreak{}常\allowbreak{}编\allowbreak{}号\allowbreak{}排\allowbreak{}除\allowbreak{}敏\allowbreak{}感\allowbreak{}性\allowbreak{}检\allowbreak{}查，\allowbreak{}简\allowbreak{}要\allowbreak{}判\allowbreak{}断\allowbreak{}其\allowbreak{}是\allowbreak{}否\allowbreak{}改\allowbreak{}变\allowbreak{}比\allowbreak{}较\allowbreak{}结\allowbreak{}论，\allowbreak{}不\allowbreak{}需\allowbreak{}扩\allowbreak{}展\allowbreak{}无\allowbreak{}关\allowbreak{}分\allowbreak{}析。\allowbreak{}不\allowbreak{}要\allowbreak{}用\allowbreak{}未\allowbreak{}明\allowbreak{}确\allowbreak{}定\allowbreak{}义\allowbreak{}的\allowbreak{}既\allowbreak{}有\allowbreak{}9项\allowbreak{}DESeq结\allowbreak{}果\allowbreak{}替\allowbreak{}代\allowbreak{}指\allowbreak{}定\allowbreak{}对\allowbreak{}比，\allowbreak{}也\allowbreak{}不\allowbreak{}要\allowbreak{}将\allowbreak{}不\allowbreak{}显\allowbreak{}著\allowbreak{}解\allowbreak{}释\allowbreak{}为\allowbreak{}生\allowbreak{}物\allowbreak{}学\allowbreak{}等\allowbreak{}效。\allowbreak{}请\allowbreak{}提\allowbreak{}交\allowbreak{}实\allowbreak{}际\allowbreak{}结\allowbreak{}果\allowbreak{}汇\allowbreak{}总\allowbreak{}和\allowbreak{}直\allowbreak{}接\allowbreak{}回\allowbreak{}答\allowbreak{}题\allowbreak{}目\allowbreak{}的\allowbreak{}简\allowbreak{}短\allowbreak{}候\allowbreak{}选\allowbreak{}答\allowbreak{}案；\allowbreak{}若\allowbreak{}二\allowbreak{}项\allowbreak{}检\allowbreak{}验\allowbreak{}不\allowbreak{}支\allowbreak{}持\allowbreak{}任\allowbreak{}何\allowbreak{}比\allowbreak{}较\allowbreak{}为\allowbreak{}无\allowbreak{}明\allowbreak{}显\allowbreak{}信\allowbreak{}号，\allowbreak{}应\allowbreak{}如\allowbreak{}实\allowbreak{}区\allowbreak{}分\allowbreak{}相\allowbreak{}对\allowbreak{}较\allowbreak{}弱\allowbreak{}效\allowbreak{}应\allowbreak{}与\allowbreak{}统\allowbreak{}计\allowbreak{}不\allowbreak{}显\allowbreak{}著。\allowbreak{}保\allowbreak{}留\allowbreak{}并\allowbreak{}提\allowbreak{}供\allowbreak{}已\allowbreak{}执\allowbreak{}行\allowbreak{}的\allowbreak{}分\allowbreak{}析\allowbreak{}脚\allowbreak{}本\allowbreak{}或\allowbreak{}笔\allowbreak{}记\allowbreak{}本\allowbreak{}路\allowbreak{}径；\allowbreak{}最\allowbreak{}终\allowbreak{}按\allowbreak{}stage、\allowbreak{}answer、\allowbreak{}payload.executed\_\allowbreak{}analysis和\allowbreak{}artifacts.analysis字\allowbreak{}段\allowbreak{}提\allowbreak{}交，\allowbreak{}只\allowbreak{}有\allowbreak{}完\allowbreak{}成\allowbreak{}运\allowbreak{}行\allowbreak{}才\allowbreak{}将\allowbreak{}executed\_\allowbreak{}analysis设\allowbreak{}为\allowbreak{}true。\allowbreak{}\par
\interactionfield{Interaction}
next\_\allowbreak{}task\par
\interactionfield{Phase}
approach\_\allowbreak{}review\par
\interactionfield{Action}
continue\par
\interactionfield{Approve Approach}
true\par
\interactionfield{Evidence}
检\allowbreak{}查\allowbreak{}报\allowbreak{}告\allowbreak{}称\allowbreak{}90个\allowbreak{}表\allowbreak{}达\allowbreak{}列\allowbreak{}与\allowbreak{}元\allowbreak{}数\allowbreak{}据\allowbreak{}完\allowbreak{}整\allowbreak{}匹\allowbreak{}配，\allowbreak{}三\allowbreak{}组\allowbreak{}织\allowbreak{}各\allowbreak{}30样\allowbreak{}本，\allowbreak{}数\allowbreak{}据\allowbreak{}为\allowbreak{}含\allowbreak{}大\allowbreak{}量\allowbreak{}小\allowbreak{}数\allowbreak{}的\allowbreak{}归\allowbreak{}一\allowbreak{}化\allowbreak{}计\allowbreak{}数；\allowbreak{}跨\allowbreak{}组\allowbreak{}织\allowbreak{}编\allowbreak{}号\allowbreak{}支\allowbreak{}持\allowbreak{}但\allowbreak{}未\allowbreak{}正\allowbreak{}式\allowbreak{}确\allowbreak{}认\allowbreak{}配\allowbreak{}对。\allowbreak{}既\allowbreak{}有\allowbreak{}9项\allowbreak{}DESeq比\allowbreak{}较\allowbreak{}缺\allowbreak{}少\allowbreak{}定\allowbreak{}义。\allowbreak{}拟\allowbreak{}议\allowbreak{}方\allowbreak{}案\allowbreak{}已\allowbreak{}明\allowbreak{}确\allowbreak{}二\allowbreak{}项\allowbreak{}检\allowbreak{}验\allowbreak{}及\allowbreak{}5\%阈\allowbreak{}值\allowbreak{}的\allowbreak{}操\allowbreak{}作\allowbreak{}含\allowbreak{}义，\allowbreak{}并\allowbreak{}正\allowbreak{}确\allowbreak{}区\allowbreak{}分\allowbreak{}名\allowbreak{}义\allowbreak{}显\allowbreak{}著、\allowbreak{}BH校\allowbreak{}正\allowbreak{}与\allowbreak{}生\allowbreak{}物\allowbreak{}学\allowbreak{}等\allowbreak{}效，\allowbreak{}足\allowbreak{}以\allowbreak{}开\allowbreak{}始\allowbreak{}执\allowbreak{}行。\allowbreak{}\par
\interactionfield{Objectives Open}
[\par
  "实\allowbreak{}际\allowbreak{}执\allowbreak{}行\allowbreak{}三\allowbreak{}项\allowbreak{}组\allowbreak{}织\allowbreak{}比\allowbreak{}较\allowbreak{}的\allowbreak{}差\allowbreak{}异\allowbreak{}表\allowbreak{}达\allowbreak{}与\allowbreak{}精\allowbreak{}确\allowbreak{}二\allowbreak{}项\allowbreak{}检\allowbreak{}验\allowbreak{}",\par
  "检\allowbreak{}查\allowbreak{}诊\allowbreak{}断\allowbreak{}及\allowbreak{}关\allowbreak{}键\allowbreak{}假\allowbreak{}设\allowbreak{}敏\allowbreak{}感\allowbreak{}性\allowbreak{}并\allowbreak{}确\allowbreak{}定\allowbreak{}题\allowbreak{}目\allowbreak{}所\allowbreak{}问\allowbreak{}比\allowbreak{}较\allowbreak{}",\par
  "提\allowbreak{}交\allowbreak{}简\allowbreak{}短\allowbreak{}答\allowbreak{}案\allowbreak{}及\allowbreak{}已\allowbreak{}执\allowbreak{}行\allowbreak{}分\allowbreak{}析\allowbreak{}文\allowbreak{}件\allowbreak{}路\allowbreak{}径\allowbreak{}"\par
]\par
\end{interactionevent}
\begin{interactionevent}{Event 28: Actor}
\interactionfield{Answer}
本\allowbreak{}轮\allowbreak{}已\allowbreak{}执\allowbreak{}行\allowbreak{}可\allowbreak{}复\allowbreak{}现\allowbreak{}输\allowbreak{}入\allowbreak{}QC，\allowbreak{}但\allowbreak{}获\allowbreak{}批\allowbreak{}的\allowbreak{}limma-\allowbreak{}trend主\allowbreak{}分\allowbreak{}析\allowbreak{}被\allowbreak{}环\allowbreak{}境\allowbreak{}阻\allowbreak{}塞，\allowbreak{}不\allowbreak{}能\allowbreak{}提\allowbreak{}供\allowbreak{}有\allowbreak{}统\allowbreak{}计\allowbreak{}结\allowbreak{}果\allowbreak{}支\allowbreak{}持\allowbreak{}的\allowbreak{}候\allowbreak{}选\allowbreak{}答\allowbreak{}案。\allowbreak{}\par
\smallskip
实\allowbreak{}际\allowbreak{}检\allowbreak{}查\allowbreak{}结\allowbreak{}果：\allowbreak{}\par
-\allowbreak{} 25,402个\allowbreak{}基\allowbreak{}因、\allowbreak{}90个\allowbreak{}样\allowbreak{}本；\allowbreak{}按\allowbreak{}“至\allowbreak{}少\allowbreak{}10个\allowbreak{}样\allowbreak{}本\allowbreak{}计\allowbreak{}数\allowbreak{}≥10”的\allowbreak{}共\allowbreak{}同\allowbreak{}规\allowbreak{}则\allowbreak{}保\allowbreak{}留\allowbreak{}17,049个\allowbreak{}基\allowbreak{}因。\allowbreak{}这\allowbreak{}是\allowbreak{}过\allowbreak{}滤\allowbreak{}后\allowbreak{}数\allowbreak{}量，\allowbreak{}尚\allowbreak{}不\allowbreak{}是\allowbreak{}有\allowbreak{}效\allowbreak{}检\allowbreak{}验\allowbreak{}数。\allowbreak{}\par
-\allowbreak{} 过\allowbreak{}滤\allowbreak{}敏\allowbreak{}感\allowbreak{}性\allowbreak{}阈\allowbreak{}值\allowbreak{}改\allowbreak{}为\allowbreak{}≥5或\allowbreak{}≥20时，\allowbreak{}分\allowbreak{}别\allowbreak{}保\allowbreak{}留\allowbreak{}18,267、\allowbreak{}15,729个\allowbreak{}基\allowbreak{}因；\allowbreak{}尚\allowbreak{}未\allowbreak{}检\allowbreak{}验\allowbreak{}这\allowbreak{}些\allowbreak{}变\allowbreak{}化\allowbreak{}是\allowbreak{}否\allowbreak{}改\allowbreak{}变\allowbreak{}结\allowbreak{}论。\allowbreak{}\par
-\allowbreak{} 2,210个\allowbreak{}基\allowbreak{}因\allowbreak{}全\allowbreak{}零，\allowbreak{}表\allowbreak{}达\allowbreak{}矩\allowbreak{}阵\allowbreak{}无\allowbreak{}非\allowbreak{}有\allowbreak{}限\allowbreak{}值；\allowbreak{}90个\allowbreak{}表\allowbreak{}达\allowbreak{}列\allowbreak{}与\allowbreak{}元\allowbreak{}数\allowbreak{}据\allowbreak{}Code顺\allowbreak{}序\allowbreak{}完\allowbreak{}全\allowbreak{}匹\allowbreak{}配。\allowbreak{}\par
-\allowbreak{} 三\allowbreak{}个\allowbreak{}组\allowbreak{}织\allowbreak{}各\allowbreak{}有\allowbreak{}30个\allowbreak{}不\allowbreak{}同\allowbreak{}的\allowbreak{}候\allowbreak{}选\allowbreak{}Subject编\allowbreak{}号，\allowbreak{}编\allowbreak{}号\allowbreak{}集\allowbreak{}合\allowbreak{}一\allowbreak{}致；\allowbreak{}同\allowbreak{}一\allowbreak{}动\allowbreak{}物\allowbreak{}配\allowbreak{}对\allowbreak{}及\allowbreak{}跨\allowbreak{}组\allowbreak{}织\allowbreak{}归\allowbreak{}一\allowbreak{}化\allowbreak{}可\allowbreak{}比\allowbreak{}性\allowbreak{}仍\allowbreak{}为\allowbreak{}假\allowbreak{}设。\allowbreak{}\par
\smallskip
环\allowbreak{}境\allowbreak{}诊\allowbreak{}断：\allowbreak{}系\allowbreak{}统\allowbreak{}无\allowbreak{}Rscript；\allowbreak{}全\allowbreak{}盘\allowbreak{}搜\allowbreak{}索\allowbreak{}未\allowbreak{}发\allowbreak{}现\allowbreak{}Rscript、\allowbreak{}limma或\allowbreak{}R安\allowbreak{}装\allowbreak{}包\allowbreak{}缓\allowbreak{}存。\allowbreak{}sudo安\allowbreak{}装\allowbreak{}被\allowbreak{}容\allowbreak{}器\allowbreak{}的\allowbreak{}no-\allowbreak{}new-\allowbreak{}privileges限\allowbreak{}制\allowbreak{}禁\allowbreak{}止；\allowbreak{}用\allowbreak{}户\allowbreak{}态\allowbreak{}安\allowbreak{}装\allowbreak{}器\allowbreak{}下\allowbreak{}载\allowbreak{}也\allowbreak{}因\allowbreak{}DNS解\allowbreak{}析\allowbreak{}失\allowbreak{}败\allowbreak{}而\allowbreak{}不\allowbreak{}可\allowbreak{}行。\allowbreak{}因\allowbreak{}此\allowbreak{}没\allowbreak{}有\allowbreak{}用\allowbreak{}其\allowbreak{}他\allowbreak{}统\allowbreak{}计\allowbreak{}方\allowbreak{}法\allowbreak{}冒\allowbreak{}充\allowbreak{}limma，\allowbreak{}也\allowbreak{}没\allowbreak{}有\allowbreak{}将\allowbreak{}归\allowbreak{}一\allowbreak{}化\allowbreak{}计\allowbreak{}数\allowbreak{}取\allowbreak{}整\allowbreak{}送\allowbreak{}入\allowbreak{}计\allowbreak{}数\allowbreak{}模\allowbreak{}型。\allowbreak{}\par
\smallskip
未\allowbreak{}完\allowbreak{}成\allowbreak{}的\allowbreak{}工\allowbreak{}作：\allowbreak{}limma(trend=\allowbreak{}TRUE)模\allowbreak{}型、\allowbreak{}三\allowbreak{}个\allowbreak{}对\allowbreak{}比\allowbreak{}的\allowbreak{}效\allowbreak{}应\allowbreak{}分\allowbreak{}布\allowbreak{}/\allowbreak{}原\allowbreak{}始\allowbreak{}p/\allowbreak{}BH结\allowbreak{}果、\allowbreak{}精\allowbreak{}确\allowbreak{}二\allowbreak{}项\allowbreak{}上\allowbreak{}尾\allowbreak{}及\allowbreak{}Holm校\allowbreak{}正，\allowbreak{}以\allowbreak{}及\allowbreak{}配\allowbreak{}对、\allowbreak{}Response异\allowbreak{}质\allowbreak{}性、\allowbreak{}异\allowbreak{}常\allowbreak{}编\allowbreak{}号\allowbreak{}排\allowbreak{}除\allowbreak{}等\allowbreak{}模\allowbreak{}型\allowbreak{}敏\allowbreak{}感\allowbreak{}性\allowbreak{}检\allowbreak{}查。\allowbreak{}不\allowbreak{}能\allowbreak{}据\allowbreak{}此\allowbreak{}判\allowbreak{}断\allowbreak{}任\allowbreak{}何\allowbreak{}比\allowbreak{}较\allowbreak{}统\allowbreak{}计\allowbreak{}不\allowbreak{}显\allowbreak{}著，\allowbreak{}或\allowbreak{}断\allowbreak{}言\allowbreak{}某\allowbreak{}项\allowbreak{}效\allowbreak{}应\allowbreak{}可\allowbreak{}忽\allowbreak{}略。\allowbreak{}\par
\smallskip
直\allowbreak{}接\allowbreak{}回\allowbreak{}答\allowbreak{}题\allowbreak{}目\allowbreak{}的\allowbreak{}候\allowbreak{}选\allowbreak{}答\allowbreak{}案：\allowbreak{}**目\allowbreak{}前\allowbreak{}无\allowbreak{}法\allowbreak{}由\allowbreak{}已\allowbreak{}执\allowbreak{}行\allowbreak{}分\allowbreak{}析\allowbreak{}确\allowbreak{}定；\allowbreak{}需\allowbreak{}要\allowbreak{}完\allowbreak{}成\allowbreak{}获\allowbreak{}批\allowbreak{}的\allowbreak{}limma模\allowbreak{}型\allowbreak{}与\allowbreak{}二\allowbreak{}项\allowbreak{}检\allowbreak{}验\allowbreak{}后\allowbreak{}作\allowbreak{}答。\allowbreak{}**\par
\smallskip
已\allowbreak{}保\allowbreak{}留\allowbreak{}并\allowbreak{}执\allowbreak{}行\allowbreak{}analysis/\allowbreak{}run\_\allowbreak{}analysis.py，\allowbreak{}输\allowbreak{}出\allowbreak{}analysis/\allowbreak{}qc\_\allowbreak{}and\_\allowbreak{}status.json；\allowbreak{}脚\allowbreak{}本\allowbreak{}在\allowbreak{}完\allowbreak{}成\allowbreak{}输\allowbreak{}入\allowbreak{}QC后\allowbreak{}明\allowbreak{}确\allowbreak{}报\allowbreak{}错\allowbreak{}退\allowbreak{}出，\allowbreak{}记\allowbreak{}录\allowbreak{}缺\allowbreak{}少\allowbreak{}R，\allowbreak{}未\allowbreak{}执\allowbreak{}行\allowbreak{}主\allowbreak{}分\allowbreak{}析。\allowbreak{}故\allowbreak{}executed\_\allowbreak{}analysis保\allowbreak{}持\allowbreak{}false。\allowbreak{}请\allowbreak{}提\allowbreak{}供\allowbreak{}可\allowbreak{}用\allowbreak{}的\allowbreak{}R与\allowbreak{}limma环\allowbreak{}境\allowbreak{}后\allowbreak{}继\allowbreak{}续\allowbreak{}原\allowbreak{}批\allowbreak{}准\allowbreak{}方\allowbreak{}法；\allowbreak{}本\allowbreak{}次\allowbreak{}不\allowbreak{}是\allowbreak{}完\allowbreak{}成\allowbreak{}的\allowbreak{}benchmark答\allowbreak{}案。\allowbreak{}\par
\interactionfield{Stage}
result\par
\interactionfield{Evidence}
实\allowbreak{}际\allowbreak{}运\allowbreak{}行\allowbreak{}python analysis/\allowbreak{}run\_\allowbreak{}analysis.py：\allowbreak{}filtered\_\allowbreak{}genes=\allowbreak{}17049；\allowbreak{}阈\allowbreak{}值\allowbreak{}5/\allowbreak{}20时\allowbreak{}18267/\allowbreak{}15729；\allowbreak{}candidate\_\allowbreak{}subject\_\allowbreak{}counts三\allowbreak{}个\allowbreak{}组\allowbreak{}织\allowbreak{}均\allowbreak{}30；\allowbreak{}codes\_\allowbreak{}match\_\allowbreak{}in\_\allowbreak{}order=\allowbreak{}true；\allowbreak{}Rscript=\allowbreak{}null。\allowbreak{}脚\allowbreak{}本\allowbreak{}以\allowbreak{}退\allowbreak{}出\allowbreak{}码\allowbreak{}1停\allowbreak{}止，\allowbreak{}明\allowbreak{}确\allowbreak{}BLOCKED。\allowbreak{}Rscript调\allowbreak{}用\allowbreak{}退\allowbreak{}出\allowbreak{}码\allowbreak{}127；\allowbreak{}sudo报\allowbreak{}告\allowbreak{}no new privileges；\allowbreak{}urllib下\allowbreak{}载\allowbreak{}返\allowbreak{}回\allowbreak{}Temporary failure in name resolution。\allowbreak{}\par
\interactionfield{Assumptions}
[\par
  "候\allowbreak{}选\allowbreak{}Subject编\allowbreak{}号\allowbreak{}对\allowbreak{}应\allowbreak{}同\allowbreak{}一\allowbreak{}动\allowbreak{}物\allowbreak{}尚\allowbreak{}未\allowbreak{}获\allowbreak{}得\allowbreak{}正\allowbreak{}式\allowbreak{}元\allowbreak{}数\allowbreak{}据\allowbreak{}确\allowbreak{}认。\allowbreak{}",\par
  "归\allowbreak{}一\allowbreak{}化\allowbreak{}计\allowbreak{}数\allowbreak{}跨\allowbreak{}组\allowbreak{}织\allowbreak{}可\allowbreak{}比\allowbreak{}性\allowbreak{}未\allowbreak{}获\allowbreak{}得\allowbreak{}方\allowbreak{}法\allowbreak{}说\allowbreak{}明\allowbreak{}确\allowbreak{}认。\allowbreak{}"\par
]\par
\interactionfield{Uncertainties}
[\par
  "所\allowbreak{}有\allowbreak{}差\allowbreak{}异\allowbreak{}表\allowbreak{}达\allowbreak{}及\allowbreak{}二\allowbreak{}项\allowbreak{}检\allowbreak{}验\allowbreak{}结\allowbreak{}果\allowbreak{}均\allowbreak{}尚\allowbreak{}未\allowbreak{}计\allowbreak{}算，\allowbreak{}不\allowbreak{}能\allowbreak{}给\allowbreak{}出\allowbreak{}比\allowbreak{}较\allowbreak{}结\allowbreak{}论。\allowbreak{}",\par
  "当\allowbreak{}前\allowbreak{}环\allowbreak{}境\allowbreak{}缺\allowbreak{}少\allowbreak{}R/\allowbreak{}limma且\allowbreak{}安\allowbreak{}装\allowbreak{}路\allowbreak{}径\allowbreak{}受\allowbreak{}限；\allowbreak{}需\allowbreak{}解\allowbreak{}决\allowbreak{}运\allowbreak{}行\allowbreak{}依\allowbreak{}赖\allowbreak{}方\allowbreak{}能\allowbreak{}按\allowbreak{}批\allowbreak{}准\allowbreak{}方\allowbreak{}法\allowbreak{}继\allowbreak{}续。\allowbreak{}"\par
]\par
\end{interactionevent}
\begin{interactionevent}{Event 29: Researcher (simulator)}
\interactionfield{Message}
目\allowbreak{}前\allowbreak{}只\allowbreak{}有\allowbreak{}QC，\allowbreak{}尚\allowbreak{}未\allowbreak{}完\allowbreak{}成\allowbreak{}题\allowbreak{}目\allowbreak{}要\allowbreak{}求\allowbreak{}的\allowbreak{}差\allowbreak{}异\allowbreak{}表\allowbreak{}达\allowbreak{}和\allowbreak{}精\allowbreak{}确\allowbreak{}二\allowbreak{}项\allowbreak{}检\allowbreak{}验，\allowbreak{}需\allowbreak{}继\allowbreak{}续。\allowbreak{}无\allowbreak{}需\allowbreak{}再\allowbreak{}尝\allowbreak{}试\allowbreak{}安\allowbreak{}装\allowbreak{}R：\allowbreak{}limma是\allowbreak{}我\allowbreak{}们\allowbreak{}提\allowbreak{}出\allowbreak{}的\allowbreak{}方\allowbreak{}案\allowbreak{}而\allowbreak{}非\allowbreak{}题\allowbreak{}目\allowbreak{}固\allowbreak{}定\allowbreak{}要\allowbreak{}求；\allowbreak{}鉴\allowbreak{}于\allowbreak{}已\allowbreak{}报\allowbreak{}告\allowbreak{}的\allowbreak{}环\allowbreak{}境\allowbreak{}限\allowbreak{}制，\allowbreak{}可\allowbreak{}明\allowbreak{}确\allowbreak{}改\allowbreak{}用\allowbreak{}现\allowbreak{}有\allowbreak{}Python工\allowbreak{}具\allowbreak{}在\allowbreak{}log2(count+1)上\allowbreak{}拟\allowbreak{}合\allowbreak{}普\allowbreak{}通\allowbreak{}最\allowbreak{}小\allowbreak{}二\allowbreak{}乘\allowbreak{}模\allowbreak{}型，\allowbreak{}以\allowbreak{}候\allowbreak{}选\allowbreak{}Subject固\allowbreak{}定\allowbreak{}效\allowbreak{}应\allowbreak{}和\allowbreak{}组\allowbreak{}织\allowbreak{}因\allowbreak{}素\allowbreak{}完\allowbreak{}成\allowbreak{}三\allowbreak{}项\allowbreak{}对\allowbreak{}比。\allowbreak{}使\allowbreak{}用\allowbreak{}残\allowbreak{}差\allowbreak{}方\allowbreak{}差、\allowbreak{}正\allowbreak{}确\allowbreak{}残\allowbreak{}差\allowbreak{}自\allowbreak{}由\allowbreak{}度\allowbreak{}及\allowbreak{}双\allowbreak{}侧\allowbreak{}t检\allowbreak{}验；\allowbreak{}明\allowbreak{}确\allowbreak{}说\allowbreak{}明\allowbreak{}这\allowbreak{}是\allowbreak{}未\allowbreak{}经\allowbreak{}经\allowbreak{}验\allowbreak{}贝\allowbreak{}叶\allowbreak{}斯\allowbreak{}收\allowbreak{}缩\allowbreak{}的\allowbreak{}替\allowbreak{}代\allowbreak{}分\allowbreak{}析，\allowbreak{}不\allowbreak{}得\allowbreak{}称\allowbreak{}为\allowbreak{}limma。\allowbreak{}保\allowbreak{}留\allowbreak{}共\allowbreak{}同\allowbreak{}过\allowbreak{}滤、\allowbreak{}比\allowbreak{}较\allowbreak{}内\allowbreak{}BH校\allowbreak{}正\allowbreak{}及\allowbreak{}既\allowbreak{}定\allowbreak{}二\allowbreak{}项\allowbreak{}检\allowbreak{}验\allowbreak{}定\allowbreak{}义\allowbreak{}（原\allowbreak{}始\allowbreak{}p<0.05为\allowbreak{}成\allowbreak{}功，\allowbreak{}p\ensuremath{{}_0}=\allowbreak{}0.05，\allowbreak{}上\allowbreak{}尾\allowbreak{}检\allowbreak{}验，\allowbreak{}另\allowbreak{}报\allowbreak{}Holm校\allowbreak{}正）。\allowbreak{}报\allowbreak{}告\allowbreak{}模\allowbreak{}型\allowbreak{}秩、\allowbreak{}残\allowbreak{}差\allowbreak{}自\allowbreak{}由\allowbreak{}度、\allowbreak{}有\allowbreak{}效\allowbreak{}检\allowbreak{}验\allowbreak{}数、\allowbreak{}三\allowbreak{}项\allowbreak{}比\allowbreak{}较\allowbreak{}的\allowbreak{}k/\allowbreak{}n、\allowbreak{}BH显\allowbreak{}著\allowbreak{}数\allowbreak{}及\allowbreak{}二\allowbreak{}项\allowbreak{}p，\allowbreak{}结\allowbreak{}合\allowbreak{}效\allowbreak{}应\allowbreak{}大\allowbreak{}小\allowbreak{}回\allowbreak{}答\allowbreak{}哪\allowbreak{}项\allowbreak{}比\allowbreak{}较\allowbreak{}效\allowbreak{}应\allowbreak{}较\allowbreak{}弱，\allowbreak{}不\allowbreak{}能\allowbreak{}将\allowbreak{}不\allowbreak{}显\allowbreak{}著\allowbreak{}直\allowbreak{}接\allowbreak{}等\allowbreak{}同\allowbreak{}等\allowbreak{}效。\allowbreak{}执\allowbreak{}行\allowbreak{}可\allowbreak{}行\allowbreak{}的\allowbreak{}基\allowbreak{}本\allowbreak{}QC和\allowbreak{}已\allowbreak{}约\allowbreak{}定\allowbreak{}敏\allowbreak{}感\allowbreak{}性\allowbreak{}检\allowbreak{}查，\allowbreak{}优\allowbreak{}先\allowbreak{}完\allowbreak{}成\allowbreak{}主\allowbreak{}分\allowbreak{}析\allowbreak{}与\allowbreak{}直\allowbreak{}接\allowbreak{}答\allowbreak{}案。\allowbreak{}保\allowbreak{}留\allowbreak{}已\allowbreak{}运\allowbreak{}行\allowbreak{}脚\allowbreak{}本，\allowbreak{}最\allowbreak{}终\allowbreak{}提\allowbreak{}交\allowbreak{}stage=\allowbreak{}final、\allowbreak{}简\allowbreak{}短\allowbreak{}answer、\allowbreak{}payload.executed\_\allowbreak{}analysis及\allowbreak{}artifacts.analysis；\allowbreak{}仅\allowbreak{}完\allowbreak{}成\allowbreak{}实\allowbreak{}际\allowbreak{}统\allowbreak{}计\allowbreak{}分\allowbreak{}析\allowbreak{}后\allowbreak{}将\allowbreak{}executed\_\allowbreak{}analysis设\allowbreak{}为\allowbreak{}true。\allowbreak{}\par
\interactionfield{Interaction}
refine\_\allowbreak{}analysis\par
\interactionfield{Phase}
result\_\allowbreak{}review\par
\interactionfield{Action}
continue\par
\interactionfield{Approve Approach}
false\par
\interactionfield{Evidence}
代\allowbreak{}理\allowbreak{}报\allowbreak{}告\allowbreak{}已\allowbreak{}执\allowbreak{}行\allowbreak{}QC并\allowbreak{}保\allowbreak{}留\allowbreak{}17,049个\allowbreak{}基\allowbreak{}因，\allowbreak{}但\allowbreak{}脚\allowbreak{}本\allowbreak{}因\allowbreak{}缺\allowbreak{}少\allowbreak{}R退\allowbreak{}出，\allowbreak{}三\allowbreak{}项\allowbreak{}差\allowbreak{}异\allowbreak{}表\allowbreak{}达\allowbreak{}及\allowbreak{}二\allowbreak{}项\allowbreak{}检\allowbreak{}验\allowbreak{}均\allowbreak{}未\allowbreak{}执\allowbreak{}行。\allowbreak{}R安\allowbreak{}装\allowbreak{}及\allowbreak{}下\allowbreak{}载\allowbreak{}受\allowbreak{}限，\allowbreak{}属\allowbreak{}于\allowbreak{}具\allowbreak{}体\allowbreak{}执\allowbreak{}行\allowbreak{}障\allowbreak{}碍；\allowbreak{}Python替\allowbreak{}代\allowbreak{}模\allowbreak{}型\allowbreak{}可\allowbreak{}保\allowbreak{}留\allowbreak{}配\allowbreak{}对\allowbreak{}设\allowbreak{}计\allowbreak{}与\allowbreak{}题\allowbreak{}目\allowbreak{}明\allowbreak{}确\allowbreak{}要\allowbreak{}求\allowbreak{}的\allowbreak{}检\allowbreak{}验，\allowbreak{}而\allowbreak{}不\allowbreak{}冒\allowbreak{}称\allowbreak{}原\allowbreak{}方\allowbreak{}案\allowbreak{}已\allowbreak{}执\allowbreak{}行。\allowbreak{}\par
\interactionfield{Objectives Open}
[\par
  "执\allowbreak{}行\allowbreak{}三\allowbreak{}项\allowbreak{}组\allowbreak{}织\allowbreak{}比\allowbreak{}较\allowbreak{}的\allowbreak{}差\allowbreak{}异\allowbreak{}表\allowbreak{}达\allowbreak{}分\allowbreak{}析\allowbreak{}及\allowbreak{}精\allowbreak{}确\allowbreak{}二\allowbreak{}项\allowbreak{}检\allowbreak{}验\allowbreak{}",\par
  "报\allowbreak{}告\allowbreak{}统\allowbreak{}计\allowbreak{}诊\allowbreak{}断\allowbreak{}和\allowbreak{}关\allowbreak{}键\allowbreak{}敏\allowbreak{}感\allowbreak{}性\allowbreak{}结\allowbreak{}果\allowbreak{}并\allowbreak{}给\allowbreak{}出\allowbreak{}直\allowbreak{}接\allowbreak{}结\allowbreak{}论\allowbreak{}",\par
  "提\allowbreak{}供\allowbreak{}符\allowbreak{}合\allowbreak{}提\allowbreak{}交\allowbreak{}字\allowbreak{}段\allowbreak{}要\allowbreak{}求\allowbreak{}的\allowbreak{}最\allowbreak{}终\allowbreak{}答\allowbreak{}案\allowbreak{}及\allowbreak{}已\allowbreak{}执\allowbreak{}行\allowbreak{}脚\allowbreak{}本\allowbreak{}路\allowbreak{}径\allowbreak{}"\par
]\par
\end{interactionevent}
\begin{interactionevent}{Event 36: Actor}
\interactionfield{Answer}
候\allowbreak{}选\allowbreak{}最\allowbreak{}终\allowbreak{}答\allowbreak{}案：\allowbreak{}**Final vs Baseline Blood（末\allowbreak{}次\allowbreak{}血\allowbreak{}液\allowbreak{}与\allowbreak{}基\allowbreak{}线\allowbreak{}血\allowbreak{}液）显\allowbreak{}示\allowbreak{}相\allowbreak{}对\allowbreak{}最\allowbreak{}弱\allowbreak{}的\allowbreak{}表\allowbreak{}达\allowbreak{}变\allowbreak{}化；\allowbreak{}但\allowbreak{}本\allowbreak{}次\allowbreak{}精\allowbreak{}确\allowbreak{}二\allowbreak{}项\allowbreak{}检\allowbreak{}验\allowbreak{}中\allowbreak{}三\allowbreak{}项\allowbreak{}比\allowbreak{}较\allowbreak{}均\allowbreak{}显\allowbreak{}著，\allowbreak{}不\allowbreak{}能\allowbreak{}将\allowbreak{}该\allowbreak{}比\allowbreak{}较\allowbreak{}称\allowbreak{}为\allowbreak{}统\allowbreak{}计\allowbreak{}不\allowbreak{}显\allowbreak{}著\allowbreak{}或\allowbreak{}生\allowbreak{}物\allowbreak{}学\allowbreak{}等\allowbreak{}效。\allowbreak{}**\par
\interactionfield{Stage}
result\par
\interactionfield{Evidence}
实\allowbreak{}际\allowbreak{}执\allowbreak{}行\allowbreak{}成\allowbreak{}功\allowbreak{}（退\allowbreak{}出\allowbreak{}码\allowbreak{}0）：\allowbreak{}analysis/\allowbreak{}ols\_\allowbreak{}analysis.py。\allowbreak{}方\allowbreak{}法\allowbreak{}为\allowbreak{}log2(count+1)上\allowbreak{}的\allowbreak{}普\allowbreak{}通\allowbreak{}OLS，\allowbreak{}非\allowbreak{}limma，\allowbreak{}无\allowbreak{}经\allowbreak{}验\allowbreak{}贝\allowbreak{}叶\allowbreak{}斯\allowbreak{}收\allowbreak{}缩；\allowbreak{}Subject固\allowbreak{}定\allowbreak{}效\allowbreak{}应\allowbreak{}加\allowbreak{}组\allowbreak{}织\allowbreak{}因\allowbreak{}素，\allowbreak{}残\allowbreak{}差\allowbreak{}方\allowbreak{}差\allowbreak{}SSE/\allowbreak{}df、\allowbreak{}对\allowbreak{}比\allowbreak{}标\allowbreak{}准\allowbreak{}误\allowbreak{}及\allowbreak{}双\allowbreak{}侧\allowbreak{}t检\allowbreak{}验。\allowbreak{}主\allowbreak{}模\allowbreak{}型\allowbreak{}90个\allowbreak{}样\allowbreak{}本，\allowbreak{}秩\allowbreak{}32，\allowbreak{}残\allowbreak{}差\allowbreak{}自\allowbreak{}由\allowbreak{}度\allowbreak{}58。\allowbreak{}共\allowbreak{}同\allowbreak{}过\allowbreak{}滤\allowbreak{}至\allowbreak{}少\allowbreak{}10个\allowbreak{}样\allowbreak{}本\allowbreak{}计\allowbreak{}数\allowbreak{}≥10保\allowbreak{}留\allowbreak{}17049个\allowbreak{}基\allowbreak{}因，\allowbreak{}三\allowbreak{}个\allowbreak{}对\allowbreak{}比\allowbreak{}有\allowbreak{}效\allowbreak{}检\allowbreak{}验\allowbreak{}数\allowbreak{}均\allowbreak{}17049。\allowbreak{}比\allowbreak{}较\allowbreak{}内\allowbreak{}BH校\allowbreak{}正\allowbreak{}逐\allowbreak{}基\allowbreak{}因\allowbreak{}结\allowbreak{}果\allowbreak{}已\allowbreak{}写\allowbreak{}入\allowbreak{}三\allowbreak{}个\allowbreak{}TSV。\allowbreak{}\par
\smallskip
主\allowbreak{}结\allowbreak{}果\allowbreak{}（依\allowbreak{}次\allowbreak{}为\allowbreak{}Final−Baseline、\allowbreak{}Dentate−Baseline、\allowbreak{}Dentate−Final）：\allowbreak{}\par
• 原\allowbreak{}始\allowbreak{}p<0.05的\allowbreak{}k/\allowbreak{}n：\allowbreak{}10689/\allowbreak{}17049=\allowbreak{}62.696\%；\allowbreak{}16059/\allowbreak{}17049=\allowbreak{}94.193\%；\allowbreak{}16066/\allowbreak{}17049=\allowbreak{}94.234\%。\allowbreak{}\par
• BH<0.05基\allowbreak{}因\allowbreak{}数：\allowbreak{}10105、\allowbreak{}16048、\allowbreak{}16049。\allowbreak{}\par
• 绝\allowbreak{}对\allowbreak{}log2表\allowbreak{}达\allowbreak{}差\allowbreak{}异\allowbreak{}中\allowbreak{}位\allowbreak{}数：\allowbreak{}0.248、\allowbreak{}1.636、\allowbreak{}1.691。\allowbreak{}\par
• 有\allowbreak{}符\allowbreak{}号\allowbreak{}效\allowbreak{}应\allowbreak{}5\%/\allowbreak{}50\%/\allowbreak{}95\%分\allowbreak{}位\allowbreak{}数：\allowbreak{}[-\allowbreak{}0.654,-\allowbreak{}0.063,0.884]、\allowbreak{}[-\allowbreak{}5.080,0.151,7.380]、\allowbreak{}[-\allowbreak{}5.177,0.194,7.500]。\allowbreak{}\par
• |log2表\allowbreak{}达\allowbreak{}差\allowbreak{}异\allowbreak{}|≥1比\allowbreak{}例：\allowbreak{}4.52\%、\allowbreak{}65.02\%、\allowbreak{}65.82\%；\allowbreak{}该\allowbreak{}描\allowbreak{}述\allowbreak{}阈\allowbreak{}值\allowbreak{}不\allowbreak{}是\allowbreak{}DEG判\allowbreak{}定\allowbreak{}门\allowbreak{}槛，\allowbreak{}也\allowbreak{}不\allowbreak{}是\allowbreak{}预\allowbreak{}设\allowbreak{}等\allowbreak{}效\allowbreak{}界\allowbreak{}限。\allowbreak{}\par
\smallskip
二\allowbreak{}项\allowbreak{}检\allowbreak{}验\allowbreak{}以\allowbreak{}有\allowbreak{}效\allowbreak{}基\allowbreak{}因\allowbreak{}为\allowbreak{}试\allowbreak{}验\allowbreak{}单\allowbreak{}位，\allowbreak{}原\allowbreak{}始\allowbreak{}双\allowbreak{}侧\allowbreak{}p<0.05为\allowbreak{}成\allowbreak{}功，\allowbreak{}H0成\allowbreak{}功\allowbreak{}概\allowbreak{}率\allowbreak{}0.05，\allowbreak{}上\allowbreak{}尾\allowbreak{}检\allowbreak{}验；\allowbreak{}Holm校\allowbreak{}正\allowbreak{}三\allowbreak{}个\allowbreak{}二\allowbreak{}项\allowbreak{}检\allowbreak{}验。\allowbreak{}各\allowbreak{}p数\allowbreak{}值\allowbreak{}极\allowbreak{}小，\allowbreak{}常\allowbreak{}规\allowbreak{}双\allowbreak{}精\allowbreak{}度\allowbreak{}输\allowbreak{}出\allowbreak{}0是\allowbreak{}下\allowbreak{}溢\allowbreak{}而\allowbreak{}非\allowbreak{}真\allowbreak{}正\allowbreak{}为\allowbreak{}零；\allowbreak{}脚\allowbreak{}本\allowbreak{}进\allowbreak{}一\allowbreak{}步\allowbreak{}用\allowbreak{}log-\allowbreak{}PMF及\allowbreak{}logsumexp累\allowbreak{}加\allowbreak{}精\allowbreak{}确\allowbreak{}二\allowbreak{}项\allowbreak{}上\allowbreak{}尾。\allowbreak{}原\allowbreak{}始\allowbreak{}log10(p)：\allowbreak{}-\allowbreak{}9159.634、\allowbreak{}-\allowbreak{}19276.254、\allowbreak{}-\allowbreak{}19293.686；\allowbreak{}Holm后\allowbreak{}log10(p)：\allowbreak{}-\allowbreak{}9159.634、\allowbreak{}-\allowbreak{}19275.953、\allowbreak{}-\allowbreak{}19293.209。\allowbreak{}全\allowbreak{}部\allowbreak{}远\allowbreak{}低\allowbreak{}于\allowbreak{}5\%阈\allowbreak{}值，\allowbreak{}故\allowbreak{}不\allowbreak{}支\allowbreak{}持\allowbreak{}任\allowbreak{}何\allowbreak{}比\allowbreak{}较\allowbreak{}“未\allowbreak{}显\allowbreak{}示\allowbreak{}超\allowbreak{}出\allowbreak{}偶\allowbreak{}然\allowbreak{}水\allowbreak{}平\allowbreak{}的\allowbreak{}差\allowbreak{}异\allowbreak{}信\allowbreak{}号\allowbreak{}”。\allowbreak{}\par
\smallskip
敏\allowbreak{}感\allowbreak{}性\allowbreak{}检\allowbreak{}查\allowbreak{}均\allowbreak{}已\allowbreak{}执\allowbreak{}行：\allowbreak{}\par
• 过\allowbreak{}滤\allowbreak{}阈\allowbreak{}值\allowbreak{}5、\allowbreak{}20分\allowbreak{}别\allowbreak{}保\allowbreak{}留\allowbreak{}18267、\allowbreak{}15729基\allowbreak{}因；\allowbreak{}血\allowbreak{}液\allowbreak{}比\allowbreak{}较\allowbreak{}名\allowbreak{}义\allowbreak{}显\allowbreak{}著\allowbreak{}比\allowbreak{}例\allowbreak{}60.71\%、\allowbreak{}64.42\%，\allowbreak{}两\allowbreak{}项\allowbreak{}脑\allowbreak{}血\allowbreak{}比\allowbreak{}较\allowbreak{}93.40\%–95.05\%；\allowbreak{}血\allowbreak{}液\allowbreak{}绝\allowbreak{}对\allowbreak{}效\allowbreak{}应\allowbreak{}中\allowbreak{}位\allowbreak{}数\allowbreak{}0.250、\allowbreak{}0.243，\allowbreak{}脑\allowbreak{}血\allowbreak{}1.614–1.698。\allowbreak{}相\allowbreak{}对\allowbreak{}排\allowbreak{}序\allowbreak{}及\allowbreak{}全\allowbreak{}部\allowbreak{}二\allowbreak{}项\allowbreak{}显\allowbreak{}著\allowbreak{}结\allowbreak{}论\allowbreak{}不\allowbreak{}变。\allowbreak{}\par
• 不\allowbreak{}配\allowbreak{}对、\allowbreak{}调\allowbreak{}整\allowbreak{}Response的\allowbreak{}模\allowbreak{}型\allowbreak{}秩\allowbreak{}5、\allowbreak{}df85：\allowbreak{}名\allowbreak{}义\allowbreak{}显\allowbreak{}著\allowbreak{}比\allowbreak{}例\allowbreak{}62.75\%、\allowbreak{}94.25\%、\allowbreak{}94.32\%，\allowbreak{}BH数\allowbreak{}10151、\allowbreak{}16058、\allowbreak{}16068；\allowbreak{}效\allowbreak{}应\allowbreak{}点\allowbreak{}估\allowbreak{}计\allowbreak{}与\allowbreak{}平\allowbreak{}衡\allowbreak{}设\allowbreak{}计\allowbreak{}的\allowbreak{}配\allowbreak{}对\allowbreak{}主\allowbreak{}模\allowbreak{}型\allowbreak{}相\allowbreak{}同，\allowbreak{}结\allowbreak{}论\allowbreak{}不\allowbreak{}变。\allowbreak{}\par
• 排\allowbreak{}除\allowbreak{}命\allowbreak{}名\allowbreak{}异\allowbreak{}常\allowbreak{}所\allowbreak{}涉\allowbreak{}候\allowbreak{}选\allowbreak{}Subject 102的\allowbreak{}全\allowbreak{}部\allowbreak{}三\allowbreak{}个\allowbreak{}样\allowbreak{}本：\allowbreak{}87样\allowbreak{}本、\allowbreak{}秩\allowbreak{}31、\allowbreak{}df56；\allowbreak{}名\allowbreak{}义\allowbreak{}显\allowbreak{}著\allowbreak{}比\allowbreak{}例\allowbreak{}62.82\%、\allowbreak{}94.24\%、\allowbreak{}94.24\%；\allowbreak{}BH数\allowbreak{}10156、\allowbreak{}16053、\allowbreak{}16055；\allowbreak{}效\allowbreak{}应\allowbreak{}中\allowbreak{}位\allowbreak{}数\allowbreak{}0.251、\allowbreak{}1.639、\allowbreak{}1.696，\allowbreak{}结\allowbreak{}论\allowbreak{}不\allowbreak{}变。\allowbreak{}\par
• Response分\allowbreak{}层\allowbreak{}配\allowbreak{}对\allowbreak{}模\allowbreak{}型\allowbreak{}每\allowbreak{}层\allowbreak{}30样\allowbreak{}本、\allowbreak{}秩\allowbreak{}12、\allowbreak{}df18。\allowbreak{}Control/\allowbreak{}bad/\allowbreak{}good层\allowbreak{}血\allowbreak{}液\allowbreak{}比\allowbreak{}较\allowbreak{}名\allowbreak{}义\allowbreak{}显\allowbreak{}著\allowbreak{}比\allowbreak{}例\allowbreak{}33.98\%/\allowbreak{}51.88\%/\allowbreak{}49.42\%，\allowbreak{}绝\allowbreak{}对\allowbreak{}效\allowbreak{}应\allowbreak{}中\allowbreak{}位\allowbreak{}数\allowbreak{}0.200/\allowbreak{}0.309/\allowbreak{}0.315；\allowbreak{}脑\allowbreak{}血\allowbreak{}比\allowbreak{}较\allowbreak{}名\allowbreak{}义\allowbreak{}显\allowbreak{}著\allowbreak{}比\allowbreak{}例\allowbreak{}89.55\%–90.71\%，\allowbreak{}绝\allowbreak{}对\allowbreak{}效\allowbreak{}应\allowbreak{}中\allowbreak{}位\allowbreak{}数\allowbreak{}1.602–1.722。\allowbreak{}各\allowbreak{}层\allowbreak{}二\allowbreak{}项\allowbreak{}检\allowbreak{}验\allowbreak{}均\allowbreak{}显\allowbreak{}著；\allowbreak{}存\allowbreak{}在\allowbreak{}组\allowbreak{}间\allowbreak{}效\allowbreak{}应\allowbreak{}强\allowbreak{}弱\allowbreak{}差\allowbreak{}异，\allowbreak{}但\allowbreak{}血\allowbreak{}液\allowbreak{}比\allowbreak{}较\allowbreak{}相\allowbreak{}对\allowbreak{}较\allowbreak{}弱\allowbreak{}的\allowbreak{}判\allowbreak{}断\allowbreak{}一\allowbreak{}致。\allowbreak{}这\allowbreak{}是\allowbreak{}分\allowbreak{}层\allowbreak{}敏\allowbreak{}感\allowbreak{}性\allowbreak{}检\allowbreak{}查，\allowbreak{}不\allowbreak{}是\allowbreak{}正\allowbreak{}式\allowbreak{}交\allowbreak{}互\allowbreak{}作\allowbreak{}用\allowbreak{}检\allowbreak{}验。\allowbreak{}\par
\smallskip
基\allowbreak{}本\allowbreak{}QC：\allowbreak{}表\allowbreak{}达\allowbreak{}值\allowbreak{}非\allowbreak{}负\allowbreak{}且\allowbreak{}无\allowbreak{}非\allowbreak{}有\allowbreak{}限\allowbreak{}值，\allowbreak{}样\allowbreak{}本\allowbreak{}与\allowbreak{}Code匹\allowbreak{}配；\allowbreak{}比\allowbreak{}对\allowbreak{}率\allowbreak{}范\allowbreak{}围\allowbreak{}60.94\%–94.68\%；\allowbreak{}样\allowbreak{}本\allowbreak{}log2表\allowbreak{}达\allowbreak{}中\allowbreak{}位\allowbreak{}数\allowbreak{}范\allowbreak{}围\allowbreak{}3.874–5.329。\allowbreak{}共\allowbreak{}同\allowbreak{}过\allowbreak{}滤\allowbreak{}矩\allowbreak{}阵\allowbreak{}中\allowbreak{}心\allowbreak{}化\allowbreak{}PCA前\allowbreak{}三\allowbreak{}主\allowbreak{}成\allowbreak{}分\allowbreak{}解\allowbreak{}释\allowbreak{}91.32\%、\allowbreak{}1.45\%、\allowbreak{}0.76\%变\allowbreak{}异，\allowbreak{}PC1主\allowbreak{}要\allowbreak{}区\allowbreak{}分\allowbreak{}脑\allowbreak{}与\allowbreak{}血，\allowbreak{}PC2上\allowbreak{}两\allowbreak{}次\allowbreak{}血\allowbreak{}液\allowbreak{}均\allowbreak{}值\allowbreak{}分\allowbreak{}离。\allowbreak{}样\allowbreak{}本\allowbreak{}级\allowbreak{}零\allowbreak{}比\allowbreak{}例、\allowbreak{}比\allowbreak{}对\allowbreak{}率、\allowbreak{}PC坐\allowbreak{}标\allowbreak{}保\allowbreak{}存\allowbreak{}在\allowbreak{}sample\_\allowbreak{}qc.tsv。\allowbreak{}没\allowbreak{}有\allowbreak{}依\allowbreak{}据\allowbreak{}表\allowbreak{}达\allowbreak{}差\allowbreak{}异\allowbreak{}事\allowbreak{}后\allowbreak{}挑\allowbreak{}选\allowbreak{}并\allowbreak{}剔\allowbreak{}除\allowbreak{}其\allowbreak{}他\allowbreak{}样\allowbreak{}本。\allowbreak{}\par
\interactionfield{Assumptions}
[\par
  "Sample ID提\allowbreak{}取\allowbreak{}的\allowbreak{}30个\allowbreak{}候\allowbreak{}选\allowbreak{}Subject代\allowbreak{}表\allowbreak{}同\allowbreak{}一\allowbreak{}对\allowbreak{}象\allowbreak{}是\allowbreak{}有\allowbreak{}编\allowbreak{}号\allowbreak{}支\allowbreak{}持、\allowbreak{}但\allowbreak{}未\allowbreak{}由\allowbreak{}正\allowbreak{}式\allowbreak{}ID字\allowbreak{}典\allowbreak{}确\allowbreak{}认\allowbreak{}的\allowbreak{}假\allowbreak{}设。\allowbreak{}",\par
  "NormCount跨\allowbreak{}组\allowbreak{}织\allowbreak{}归\allowbreak{}一\allowbreak{}化\allowbreak{}尺\allowbreak{}度\allowbreak{}可\allowbreak{}比\allowbreak{}是\allowbreak{}未\allowbreak{}确\allowbreak{}认\allowbreak{}假\allowbreak{}设；\allowbreak{}效\allowbreak{}应\allowbreak{}为\allowbreak{}log2(count+1)均\allowbreak{}值\allowbreak{}差，\allowbreak{}而\allowbreak{}非\allowbreak{}原\allowbreak{}始\allowbreak{}计\allowbreak{}数\allowbreak{}模\allowbreak{}型\allowbreak{}中\allowbreak{}的\allowbreak{}精\allowbreak{}确\allowbreak{}log2 fold-\allowbreak{}change。\allowbreak{}",\par
  "普\allowbreak{}通\allowbreak{}OLS t检\allowbreak{}验\allowbreak{}依\allowbreak{}赖\allowbreak{}残\allowbreak{}差\allowbreak{}模\allowbreak{}型\allowbreak{}及\allowbreak{}同\allowbreak{}方\allowbreak{}差\allowbreak{}等\allowbreak{}假\allowbreak{}设；\allowbreak{}未\allowbreak{}作\allowbreak{}limma式\allowbreak{}均\allowbreak{}值\allowbreak{}—方\allowbreak{}差\allowbreak{}趋\allowbreak{}势\allowbreak{}建\allowbreak{}模\allowbreak{}或\allowbreak{}经\allowbreak{}验\allowbreak{}贝\allowbreak{}叶\allowbreak{}斯\allowbreak{}收\allowbreak{}缩。\allowbreak{}",\par
  "基\allowbreak{}因\allowbreak{}相\allowbreak{}关\allowbreak{}性\allowbreak{}违\allowbreak{}反\allowbreak{}二\allowbreak{}项\allowbreak{}独\allowbreak{}立\allowbreak{}性，\allowbreak{}所\allowbreak{}谓\allowbreak{}精\allowbreak{}确\allowbreak{}p仅\allowbreak{}在\allowbreak{}二\allowbreak{}项\allowbreak{}模\allowbreak{}型\allowbreak{}下\allowbreak{}精\allowbreak{}确；\allowbreak{}5\%成\allowbreak{}功\allowbreak{}概\allowbreak{}率\allowbreak{}及\allowbreak{}原\allowbreak{}始\allowbreak{}p筛\allowbreak{}选\allowbreak{}是\allowbreak{}本\allowbreak{}次\allowbreak{}操\allowbreak{}作\allowbreak{}性\allowbreak{}定\allowbreak{}义，\allowbreak{}不\allowbreak{}代\allowbreak{}表\allowbreak{}恢\allowbreak{}复\allowbreak{}了\allowbreak{}原\allowbreak{}作\allowbreak{}者\allowbreak{}方\allowbreak{}法。\allowbreak{}"\par
]\par
\interactionfield{Uncertainties}
[\par
  "题\allowbreak{}目\allowbreak{}所\allowbreak{}称\allowbreak{}negligible不\allowbreak{}能\allowbreak{}由\allowbreak{}本\allowbreak{}次\allowbreak{}显\allowbreak{}著\allowbreak{}性\allowbreak{}结\allowbreak{}果\allowbreak{}严\allowbreak{}格\allowbreak{}支\allowbreak{}持；\allowbreak{}只\allowbreak{}能\allowbreak{}确\allowbreak{}认\allowbreak{}Final vs Baseline Blood相\allowbreak{}对\allowbreak{}更\allowbreak{}弱。\allowbreak{}",\par
  "归\allowbreak{}一\allowbreak{}化、\allowbreak{}批\allowbreak{}次、\allowbreak{}细\allowbreak{}胞\allowbreak{}组\allowbreak{}成\allowbreak{}及\allowbreak{}处\allowbreak{}理\allowbreak{}时\allowbreak{}间\allowbreak{}等\allowbreak{}信\allowbreak{}息\allowbreak{}不\allowbreak{}足，\allowbreak{}可\allowbreak{}混\allowbreak{}杂\allowbreak{}组\allowbreak{}织\allowbreak{}或\allowbreak{}时\allowbreak{}间\allowbreak{}差\allowbreak{}异；\allowbreak{}主\allowbreak{}模\allowbreak{}型\allowbreak{}按\allowbreak{}基\allowbreak{}因\allowbreak{}在\allowbreak{}三\allowbreak{}组\allowbreak{}织\allowbreak{}间\allowbreak{}合\allowbreak{}并\allowbreak{}估\allowbreak{}计\allowbreak{}残\allowbreak{}差\allowbreak{}方\allowbreak{}差，\allowbreak{}未\allowbreak{}建\allowbreak{}模\allowbreak{}组\allowbreak{}织\allowbreak{}特\allowbreak{}异\allowbreak{}方\allowbreak{}差。\allowbreak{}",\par
  "Response分\allowbreak{}层\allowbreak{}提\allowbreak{}示\allowbreak{}异\allowbreak{}质\allowbreak{}性，\allowbreak{}但\allowbreak{}本\allowbreak{}轮\allowbreak{}未\allowbreak{}执\allowbreak{}行\allowbreak{}正\allowbreak{}式\allowbreak{}交\allowbreak{}互\allowbreak{}作\allowbreak{}用\allowbreak{}检\allowbreak{}验；\allowbreak{}已\allowbreak{}有\allowbreak{}九\allowbreak{}项\allowbreak{}DESeq结\allowbreak{}果\allowbreak{}未\allowbreak{}用\allowbreak{}于\allowbreak{}替\allowbreak{}代\allowbreak{}指\allowbreak{}定\allowbreak{}比\allowbreak{}较。\allowbreak{}"\par
]\par
\end{interactionevent}
\begin{interactionevent}{Event 37: Researcher (simulator)}
\interactionfield{Message}
分\allowbreak{}析\allowbreak{}已\allowbreak{}充\allowbreak{}分\allowbreak{}报\allowbreak{}告，\allowbreak{}可\allowbreak{}以\allowbreak{}采\allowbreak{}用\allowbreak{}已\allowbreak{}提\allowbreak{}供\allowbreak{}的\allowbreak{}final\_\allowbreak{}candidate作\allowbreak{}为\allowbreak{}最\allowbreak{}终\allowbreak{}提\allowbreak{}交，\allowbreak{}保\allowbreak{}留\allowbreak{}analysis/\allowbreak{}ols\_\allowbreak{}analysis.py及\allowbreak{}executed\_\allowbreak{}analysis=\allowbreak{}true。\allowbreak{}答\allowbreak{}案\allowbreak{}应\allowbreak{}明\allowbreak{}确\allowbreak{}指\allowbreak{}出\allowbreak{}Final vs Baseline Blood相\allowbreak{}对\allowbreak{}效\allowbreak{}应\allowbreak{}最\allowbreak{}弱，\allowbreak{}同\allowbreak{}时\allowbreak{}保\allowbreak{}留\allowbreak{}本\allowbreak{}次\allowbreak{}操\allowbreak{}作\allowbreak{}性\allowbreak{}二\allowbreak{}项\allowbreak{}检\allowbreak{}验\allowbreak{}三\allowbreak{}项\allowbreak{}均\allowbreak{}显\allowbreak{}著\allowbreak{}的\allowbreak{}限\allowbreak{}定；\allowbreak{}不\allowbreak{}要\allowbreak{}将\allowbreak{}其\allowbreak{}改\allowbreak{}写\allowbreak{}为\allowbreak{}已\allowbreak{}证\allowbreak{}明\allowbreak{}无\allowbreak{}效\allowbreak{}应\allowbreak{}或\allowbreak{}生\allowbreak{}物\allowbreak{}学\allowbreak{}等\allowbreak{}效。\allowbreak{}无\allowbreak{}需\allowbreak{}追\allowbreak{}加\allowbreak{}分\allowbreak{}析。\allowbreak{}\par
\interactionfield{Interaction}
conclude\par
\interactionfield{Phase}
result\_\allowbreak{}review\par
\interactionfield{Action}
stop\par
\interactionfield{Approve Approach}
false\par
\interactionfield{Evidence}
代\allowbreak{}理\allowbreak{}报\allowbreak{}告\allowbreak{}OLS脚\allowbreak{}本\allowbreak{}成\allowbreak{}功\allowbreak{}运\allowbreak{}行，\allowbreak{}主\allowbreak{}模\allowbreak{}型\allowbreak{}秩\allowbreak{}32、\allowbreak{}残\allowbreak{}差\allowbreak{}自\allowbreak{}由\allowbreak{}度\allowbreak{}58，\allowbreak{}三\allowbreak{}项\allowbreak{}比\allowbreak{}较\allowbreak{}各\allowbreak{}检\allowbreak{}验\allowbreak{}17049个\allowbreak{}基\allowbreak{}因，\allowbreak{}并\allowbreak{}提\allowbreak{}供\allowbreak{}BH、\allowbreak{}精\allowbreak{}确\allowbreak{}二\allowbreak{}项\allowbreak{}上\allowbreak{}尾\allowbreak{}及\allowbreak{}Holm结\allowbreak{}果。\allowbreak{}血\allowbreak{}液\allowbreak{}比\allowbreak{}较\allowbreak{}绝\allowbreak{}对\allowbreak{}log2差\allowbreak{}异\allowbreak{}中\allowbreak{}位\allowbreak{}数\allowbreak{}0.248，\allowbreak{}低\allowbreak{}于\allowbreak{}两\allowbreak{}项\allowbreak{}脑\allowbreak{}血\allowbreak{}比\allowbreak{}较\allowbreak{}的\allowbreak{}1.636和\allowbreak{}1.691；\allowbreak{}所\allowbreak{}有\allowbreak{}二\allowbreak{}项\allowbreak{}检\allowbreak{}验\allowbreak{}均\allowbreak{}显\allowbreak{}著，\allowbreak{}敏\allowbreak{}感\allowbreak{}性\allowbreak{}检\allowbreak{}查\allowbreak{}保\allowbreak{}持\allowbreak{}该\allowbreak{}结\allowbreak{}论。\allowbreak{}最\allowbreak{}终\allowbreak{}候\allowbreak{}选\allowbreak{}已\allowbreak{}具\allowbreak{}备\allowbreak{}stage、\allowbreak{}answer、\allowbreak{}执\allowbreak{}行\allowbreak{}标\allowbreak{}记\allowbreak{}和\allowbreak{}分\allowbreak{}析\allowbreak{}路\allowbreak{}径，\allowbreak{}且\allowbreak{}如\allowbreak{}实\allowbreak{}区\allowbreak{}分\allowbreak{}相\allowbreak{}对\allowbreak{}较\allowbreak{}弱\allowbreak{}与\allowbreak{}统\allowbreak{}计\allowbreak{}不\allowbreak{}显\allowbreak{}著。\allowbreak{}上\allowbreak{}述\allowbreak{}为\allowbreak{}已\allowbreak{}报\allowbreak{}告\allowbreak{}执\allowbreak{}行\allowbreak{}证\allowbreak{}据，\allowbreak{}未\allowbreak{}独\allowbreak{}立\allowbreak{}复\allowbreak{}核。\allowbreak{}\par
\interactionfield{Objectives Open}
[]\par
\end{interactionevent}
\Needspace{10\baselineskip}
\subsection{Incremental Clarification: SciConvBench with Persona 155}
\label{app:case-clarification}
Task: \texttt{solMech/\allowbreak{}disambiguation/\allowbreak{}case\_\allowbreak{}080}. Trial 1; original endpoint; persona 155.
% Snapshot SHA256: a27ef0a21f4064b424c3fbe21606a182be46866aaa949d7b0b6a6337620a9a52
\begin{interactionevent}{Event 1: Initial task}
\interactionfield{Message}
Assume small-\allowbreak{}deformation, linear-\allowbreak{}elastic solid mechanics unless otherwise required.\par
\smallskip
A solid circular shaft is subjected to a torque and shear loading.\par
\smallskip
Determine the relevant response quantities for the case, focusing on: design a solid shaft diameter from an angle-\allowbreak{}of-\allowbreak{}twist limit.\par
\smallskip
State any additional assumptions you need and present the governing relations clearly.\par
\end{interactionevent}
\begin{interactionevent}{Event 2: Actor clarification question}
\interactionfield{Question}
What torque acts on the shaft (specify its magnitude, or indicate that it should remain a symbolic parameter)?\par
\end{interactionevent}
\begin{interactionevent}{Event 3: Researcher (simulator) reply}
\interactionfield{Answer}
轴\allowbreak{}上\allowbreak{}施\allowbreak{}加\allowbreak{}的\allowbreak{}扭\allowbreak{}矩\allowbreak{}为\allowbreak{} \textbackslash{}(T =\allowbreak{} 1\{,\}000\{,\}000\textbackslash{} \textbackslash{}text\{N·mm\}\textbackslash{})，\allowbreak{}请\allowbreak{}使\allowbreak{}用\allowbreak{}此\allowbreak{}数\allowbreak{}值\allowbreak{}计\allowbreak{}算，\allowbreak{}无\allowbreak{}需\allowbreak{}保\allowbreak{}留\allowbreak{}为\allowbreak{}符\allowbreak{}号\allowbreak{}参\allowbreak{}数。\allowbreak{}\par
\end{interactionevent}
\begin{interactionevent}{Event 4: Actor clarification question}
\interactionfield{Question}
轴\allowbreak{}的\allowbreak{}受\allowbreak{}扭\allowbreak{}长\allowbreak{}度\allowbreak{}是\allowbreak{}多\allowbreak{}少\allowbreak{}（单\allowbreak{}位\allowbreak{}为\allowbreak{} mm）？\allowbreak{}\par
\end{interactionevent}
\begin{interactionevent}{Event 5: Researcher (simulator) reply}
\interactionfield{Answer}
轴\allowbreak{}的\allowbreak{}受\allowbreak{}扭\allowbreak{}长\allowbreak{}度\allowbreak{}为\allowbreak{} \textbackslash{}(L =\allowbreak{} 1500\textbackslash{} \textbackslash{}text\{mm\}\textbackslash{})。\allowbreak{}\par
\end{interactionevent}
\begin{interactionevent}{Event 6: Actor clarification question}
\interactionfield{Question}
轴\allowbreak{}允\allowbreak{}许\allowbreak{}的\allowbreak{}最\allowbreak{}大\allowbreak{}扭\allowbreak{}转\allowbreak{}角\allowbreak{}是\allowbreak{}多\allowbreak{}少\allowbreak{}（请\allowbreak{}注\allowbreak{}明\allowbreak{}单\allowbreak{}位\allowbreak{}为\allowbreak{}度\allowbreak{}或\allowbreak{}弧\allowbreak{}度）？\allowbreak{}\par
\end{interactionevent}
\begin{interactionevent}{Event 7: Researcher (simulator) reply}
\interactionfield{Answer}
轴\allowbreak{}在\allowbreak{}全\allowbreak{}长\allowbreak{} \textbackslash{}(L =\allowbreak{} 1500\textbackslash{} \textbackslash{}text\{mm\}\textbackslash{}) 上\allowbreak{}允\allowbreak{}许\allowbreak{}的\allowbreak{}最\allowbreak{}大\allowbreak{}扭\allowbreak{}转\allowbreak{}角\allowbreak{}为\allowbreak{} \textbackslash{}(\textbackslash{}varphi\_\allowbreak{}\{\textbackslash{}text\{allow\}\} =\allowbreak{} 0.75\textasciicircum{}\textbackslash{}circ\textbackslash{})。\allowbreak{}代\allowbreak{}入\allowbreak{}扭\allowbreak{}转\allowbreak{}公\allowbreak{}式\allowbreak{}时，\allowbreak{}请\allowbreak{}换\allowbreak{}算\allowbreak{}为\allowbreak{}弧\allowbreak{}度，\allowbreak{}即\allowbreak{}约\allowbreak{} \textbackslash{}(0.01308997\textbackslash{} \textbackslash{}text\{rad\}\textbackslash{})。\allowbreak{}\par
\end{interactionevent}
\begin{interactionevent}{Event 8: Actor clarification question}
\interactionfield{Question}
轴\allowbreak{}材\allowbreak{}料\allowbreak{}的\allowbreak{}剪\allowbreak{}切\allowbreak{}模\allowbreak{}量\allowbreak{} \textbackslash{}(G\textbackslash{}) 是\allowbreak{}多\allowbreak{}少\allowbreak{}（单\allowbreak{}位\allowbreak{}为\allowbreak{} MPa 或\allowbreak{} N/\allowbreak{}mm²）？\allowbreak{}\par
\end{interactionevent}
\begin{interactionevent}{Event 9: Researcher (simulator) reply}
\interactionfield{Answer}
轴\allowbreak{}的\allowbreak{}材\allowbreak{}料\allowbreak{}为\allowbreak{} **6061-\allowbreak{}T6 铝\allowbreak{}合\allowbreak{}金\allowbreak{}**，\allowbreak{}剪\allowbreak{}切\allowbreak{}模\allowbreak{}量\allowbreak{}为\allowbreak{} \textbackslash{}(G =\allowbreak{} 25939.8\textbackslash{} \textbackslash{}text\{MPa\}\textbackslash{})，\allowbreak{}即\allowbreak{} \textbackslash{}(25939.8\textbackslash{} \textbackslash{}text\{N/\allowbreak{}mm\}\textasciicircum{}2\textbackslash{})。\allowbreak{}请\allowbreak{}使\allowbreak{}用\allowbreak{}此\allowbreak{}数\allowbreak{}值\allowbreak{}计\allowbreak{}算。\allowbreak{}\par
\end{interactionevent}
\begin{interactionevent}{Event 10: Actor clarification question}
\interactionfield{Question}
题\allowbreak{}目\allowbreak{}中\allowbreak{}的\allowbreak{}“剪\allowbreak{}切\allowbreak{}载\allowbreak{}荷\allowbreak{}”是\allowbreak{}指\allowbreak{}扭\allowbreak{}矩\allowbreak{}产\allowbreak{}生\allowbreak{}的\allowbreak{}剪\allowbreak{}应\allowbreak{}力，\allowbreak{}还\allowbreak{}是\allowbreak{}另\allowbreak{}有\allowbreak{}独\allowbreak{}立\allowbreak{}施\allowbreak{}加\allowbreak{}的\allowbreak{}横\allowbreak{}向\allowbreak{}剪\allowbreak{}力？\allowbreak{}\par
\end{interactionevent}
\begin{interactionevent}{Event 11: Researcher (simulator) reply}
\interactionfield{Answer}
这\allowbreak{}里\allowbreak{}的\allowbreak{}“剪\allowbreak{}切\allowbreak{}载\allowbreak{}荷\allowbreak{}”指\allowbreak{}扭\allowbreak{}矩\allowbreak{}产\allowbreak{}生\allowbreak{}的\allowbreak{}剪\allowbreak{}应\allowbreak{}力，\allowbreak{}**没\allowbreak{}有\allowbreak{}另\allowbreak{}行\allowbreak{}施\allowbreak{}加\allowbreak{}的\allowbreak{}横\allowbreak{}向\allowbreak{}剪\allowbreak{}力\allowbreak{}**。\allowbreak{}\par
\smallskip
本\allowbreak{}题\allowbreak{}只\allowbreak{}需\allowbreak{}根\allowbreak{}据\allowbreak{}扭\allowbreak{}转\allowbreak{}角\allowbreak{}限\allowbreak{}制\allowbreak{}确\allowbreak{}定\allowbreak{}实\allowbreak{}心\allowbreak{}圆\allowbreak{}轴\allowbreak{}的\allowbreak{}最\allowbreak{}小\allowbreak{}直\allowbreak{}径，\allowbreak{}使\allowbreak{}用\allowbreak{}\par
\textbackslash{}[\par
\textbackslash{}varphi=\allowbreak{}\textbackslash{}frac\{TL\}\{JG\},\textbackslash{}qquad J=\allowbreak{}\textbackslash{}frac\{\textbackslash{}pi d\textasciicircum{}4\}\{32\}.\par
\textbackslash{}]\par
不\allowbreak{}需\allowbreak{}要\allowbreak{}增\allowbreak{}加\allowbreak{}横\allowbreak{}向\allowbreak{}剪\allowbreak{}力\allowbreak{}分\allowbreak{}析\allowbreak{}或\allowbreak{}剪\allowbreak{}应\allowbreak{}力\allowbreak{}强\allowbreak{}度\allowbreak{}校\allowbreak{}核\allowbreak{}条\allowbreak{}件。\allowbreak{}\par
\end{interactionevent}
\begin{interactionevent}{Event 12: Actor final specification}
\interactionfield{Answer}
按\allowbreak{}小\allowbreak{}变\allowbreak{}形、\allowbreak{}线\allowbreak{}弹\allowbreak{}性\allowbreak{}固\allowbreak{}体\allowbreak{}力\allowbreak{}学\allowbreak{}理\allowbreak{}论，\allowbreak{}根\allowbreak{}据\allowbreak{}全\allowbreak{}长\allowbreak{}扭\allowbreak{}转\allowbreak{}角\allowbreak{}限\allowbreak{}制\allowbreak{}计\allowbreak{}算\allowbreak{}实\allowbreak{}心\allowbreak{}圆\allowbreak{}轴\allowbreak{}的\allowbreak{}最\allowbreak{}小\allowbreak{}直\allowbreak{}径，\allowbreak{}以\allowbreak{}中\allowbreak{}文\allowbreak{}呈\allowbreak{}现\allowbreak{}计\allowbreak{}算\allowbreak{}过\allowbreak{}程\allowbreak{}与\allowbreak{}结\allowbreak{}果。\allowbreak{}\par
假\allowbreak{}设\allowbreak{}轴\allowbreak{}为\allowbreak{}等\allowbreak{}直\allowbreak{}径、\allowbreak{}均\allowbreak{}质、\allowbreak{}各\allowbreak{}向\allowbreak{}同\allowbreak{}性\allowbreak{}的\allowbreak{}实\allowbreak{}心\allowbreak{}圆\allowbreak{}轴，\allowbreak{}沿\allowbreak{}全\allowbreak{}长\allowbreak{}承\allowbreak{}受\allowbreak{}恒\allowbreak{}定\allowbreak{}内\allowbreak{}扭\allowbreak{}矩，\allowbreak{}采\allowbreak{}用\allowbreak{}圣\allowbreak{}维\allowbreak{}南\allowbreak{}扭\allowbreak{}转\allowbreak{}理\allowbreak{}论\allowbreak{}并\allowbreak{}忽\allowbreak{}略\allowbreak{}端\allowbreak{}部\allowbreak{}局\allowbreak{}部\allowbreak{}效\allowbreak{}应。\allowbreak{}\par
采\allowbreak{}用\allowbreak{}扭\allowbreak{}矩\allowbreak{} T =\allowbreak{} 1,000,000 N·mm、\allowbreak{}受\allowbreak{}扭\allowbreak{}长\allowbreak{}度\allowbreak{} L =\allowbreak{} 1500 mm。\allowbreak{}\par
材\allowbreak{}料\allowbreak{}为\allowbreak{} 6061-\allowbreak{}T6 铝\allowbreak{}合\allowbreak{}金，\allowbreak{}采\allowbreak{}用\allowbreak{}剪\allowbreak{}切\allowbreak{}模\allowbreak{}量\allowbreak{} G =\allowbreak{} 25939.8 MPa =\allowbreak{} 25939.8 N/\allowbreak{}mm²。\allowbreak{}\par
全\allowbreak{}长\allowbreak{}允\allowbreak{}许\allowbreak{}扭\allowbreak{}转\allowbreak{}角\allowbreak{}为\allowbreak{} φ\_\allowbreak{}allow =\allowbreak{} 0.75°，\allowbreak{}计\allowbreak{}算\allowbreak{}时\allowbreak{}使\allowbreak{}用\allowbreak{} φ\_\allowbreak{}allow =\allowbreak{} 0.75π/\allowbreak{}180 rad ≈ 0.01308997 rad。\allowbreak{}\par
明\allowbreak{}确\allowbreak{}题\allowbreak{}目\allowbreak{}中\allowbreak{}的\allowbreak{}剪\allowbreak{}切\allowbreak{}载\allowbreak{}荷\allowbreak{}仅\allowbreak{}指\allowbreak{}扭\allowbreak{}矩\allowbreak{}产\allowbreak{}生\allowbreak{}的\allowbreak{}剪\allowbreak{}应\allowbreak{}力，\allowbreak{}不\allowbreak{}存\allowbreak{}在\allowbreak{}独\allowbreak{}立\allowbreak{}施\allowbreak{}加\allowbreak{}的\allowbreak{}横\allowbreak{}向\allowbreak{}剪\allowbreak{}力。\allowbreak{}\par
列\allowbreak{}出\allowbreak{}控\allowbreak{}制\allowbreak{}关\allowbreak{}系\allowbreak{} φ =\allowbreak{} TL/\allowbreak{}(JG)、\allowbreak{}J =\allowbreak{} πd⁴/\allowbreak{}32 以\allowbreak{}及\allowbreak{}设\allowbreak{}计\allowbreak{}约\allowbreak{}束\allowbreak{} φ ≤ φ\_\allowbreak{}allow。\allowbreak{}\par
由\allowbreak{}设\allowbreak{}计\allowbreak{}约\allowbreak{}束\allowbreak{}推\allowbreak{}导\allowbreak{}最\allowbreak{}小\allowbreak{}直\allowbreak{}径\allowbreak{}公\allowbreak{}式\allowbreak{} d\_\allowbreak{}min =\allowbreak{} [32TL/\allowbreak{}(πGφ\_\allowbreak{}allow)]\textasciicircum{}(1/\allowbreak{}4)，\allowbreak{}代\allowbreak{}入\allowbreak{}给\allowbreak{}定\allowbreak{}数\allowbreak{}值\allowbreak{}并\allowbreak{}以\allowbreak{} mm 报\allowbreak{}告\allowbreak{}结\allowbreak{}果。\allowbreak{}\par
若\allowbreak{}对\allowbreak{}计\allowbreak{}算\allowbreak{}直\allowbreak{}径\allowbreak{}取\allowbreak{}整，\allowbreak{}应\allowbreak{}向\allowbreak{}上\allowbreak{}取\allowbreak{}整，\allowbreak{}以\allowbreak{}保\allowbreak{}证\allowbreak{}满\allowbreak{}足\allowbreak{}扭\allowbreak{}转\allowbreak{}角\allowbreak{}限\allowbreak{}制。\allowbreak{}\par
不\allowbreak{}增\allowbreak{}加\allowbreak{}横\allowbreak{}向\allowbreak{}剪\allowbreak{}力\allowbreak{}分\allowbreak{}析、\allowbreak{}剪\allowbreak{}应\allowbreak{}力\allowbreak{}强\allowbreak{}度\allowbreak{}校\allowbreak{}核\allowbreak{}或\allowbreak{}额\allowbreak{}外\allowbreak{}安\allowbreak{}全\allowbreak{}系\allowbreak{}数，\allowbreak{}并\allowbreak{}说\allowbreak{}明\allowbreak{}所\allowbreak{}得\allowbreak{}最\allowbreak{}小\allowbreak{}直\allowbreak{}径\allowbreak{}仅\allowbreak{}由\allowbreak{}扭\allowbreak{}转\allowbreak{}刚\allowbreak{}度\allowbreak{}条\allowbreak{}件\allowbreak{}确\allowbreak{}定。\allowbreak{}\par
\end{interactionevent}
\Needspace{10\baselineskip}
\subsection{Matched SciConvBench Conversation without a Persona}
\label{app:case-clarification-off}
Task: \texttt{solMech/\allowbreak{}disambiguation/\allowbreak{}case\_\allowbreak{}080}. Trial 1; original endpoint; persona persona\_\allowbreak{}off.
% Snapshot SHA256: f8ff1ee9fb7bae30e4c66ac099a422d3c86c385c2c4c9561c735e657e0f42d95
\begin{interactionevent}{Event 1: Initial task}
\interactionfield{Message}
Assume small-\allowbreak{}deformation, linear-\allowbreak{}elastic solid mechanics unless otherwise required.\par
\smallskip
A solid circular shaft is subjected to a torque and shear loading.\par
\smallskip
Determine the relevant response quantities for the case, focusing on: design a solid shaft diameter from an angle-\allowbreak{}of-\allowbreak{}twist limit.\par
\smallskip
State any additional assumptions you need and present the governing relations clearly.\par
\end{interactionevent}
\begin{interactionevent}{Event 2: Actor clarification question}
\interactionfield{Question}
Should the shaft diameter be derived symbolically, or calculated numerically from specified loading, geometry, material properties, and twist limit?\par
\end{interactionevent}
\begin{interactionevent}{Event 3: Researcher (simulator) reply}
\interactionfield{Answer}
Please calculate the diameter numerically using N–mm–MPa units:\par
\smallskip
-\allowbreak{} Length: \textbackslash{}(L =\allowbreak{} 1500\textbackslash{} \textbackslash{}text\{mm\}\textbackslash{})\par
-\allowbreak{} Material: 6061-\allowbreak{}T6 aluminum\par
-\allowbreak{} Shear modulus: \textbackslash{}(G =\allowbreak{} 25939.8\textbackslash{} \textbackslash{}text\{MPa\}\textbackslash{})\par
-\allowbreak{} Torque: \textbackslash{}(T =\allowbreak{} 1\{,\}000\{,\}000\textbackslash{} \textbackslash{}text\{N·mm\}\textbackslash{})\par
-\allowbreak{} Allowable twist: \textbackslash{}(\textbackslash{}phi\_\allowbreak{}\{\textbackslash{}text\{allow\}\} =\allowbreak{} 0.75\textasciicircum{}\textbackslash{}circ\textbackslash{}) over the full length\par
\smallskip
Determine the minimum solid-\allowbreak{}shaft diameter using \textbackslash{}(\textbackslash{}phi =\allowbreak{} TL/\allowbreak{}(JG)\textbackslash{}) and \textbackslash{}(J =\allowbreak{} \textbackslash{}pi d\textasciicircum{}4/\allowbreak{}32\textbackslash{}), converting the twist limit to radians. Please also show the governing expression for the diameter.\par
\end{interactionevent}
\begin{interactionevent}{Event 4: Actor final specification}
\interactionfield{Answer}
Determine the minimum solid circular shaft diameter numerically using N–mm–MPa units.\par
Assume small-\allowbreak{}deformation, linear-\allowbreak{}elastic behavior for a homogeneous, isotropic 6061-\allowbreak{}T6 aluminum shaft with constant circular cross section.\par
Use a shaft length L =\allowbreak{} 1500 mm, shear modulus G =\allowbreak{} 25939.8 MPa, and constant internal torque T =\allowbreak{} 1,000,000 N·mm over the full length.\par
Limit the relative angle of twist over the full shaft length to 0.75 degrees.\par
Convert the allowable twist to radians using φ\_\allowbreak{}allow =\allowbreak{} 0.75π/\allowbreak{}180.\par
Use the governing torsion relations φ =\allowbreak{} TL/\allowbreak{}(JG) and J =\allowbreak{} πd⁴/\allowbreak{}32.\par
Show that the twist constraint gives d ≥ [32TL/\allowbreak{}(πGφ\_\allowbreak{}allow)]\textasciicircum{}(1/\allowbreak{}4).\par
Substitute the numerical values and report the minimum diameter in millimeters without rounding downward.\par
Verify that the reported diameter satisfies the allowable twist.\par
Treat this as sizing solely for torsional stiffness; no separate transverse shear load is specified.\par
State that strength, bending, transverse shear, stress concentrations, and connection effects are outside this sizing calculation.\par
\end{interactionevent}
